% Paper 28: Periods of Dirac Spectral Sums on S^3
% GeoVac Project.  Narrowed 2026-11-07 (PI-approved) to the number theory of
% Camporesi--Higuchi spectral sums; cut material in
% docs/chronicles/paper28_narrowing_cut.md.

\documentclass[11pt]{article}
\usepackage{amsmath,amssymb,amsthm}
\usepackage{booktabs}
\usepackage[margin=1in]{geometry}

\newtheorem{theorem}{Theorem}
\newtheorem{proposition}{Proposition}
\newtheorem{corollary}{Corollary}
\newtheorem{observation}{Observation}
\newtheorem{remark}{Remark}

% Shorthand macros for number rings and motivic categories
\newcommand{\MT}{\mathcal{MT}}
\newcommand{\Z}{\mathbb{Z}}

\title{Periods of Dirac Spectral Sums on $S^3$}
\author{GeoVac Project}
\date{April 2026; narrowed November 2026}

\begin{document}
\maketitle

\begin{abstract}
We determine the periods of spectral sums built from the
Camporesi--Higuchi Dirac spectrum on the unit three-sphere,
$|\lambda_n| = n + 3/2$ with degeneracy $2(n+1)(n+2)$.
(1)~\textbf{T9:} the spectral zeta function of $D^2$ is a two-term
polynomial in $\pi^2$ with rational coefficients at every integer
$s \ge 1$; at half-integer $s \ge 5/2$ the same closed form carries odd
zeta values (Theorem~\ref{thm:T9}).
(2)~\textbf{Parity:} the Dirac Dirichlet series $D(s)$ lies in
$\mathbb{Q}[\pi^2]$ at even $s$ and is a rational combination of
$\zeta(s{-}2)$ and $\zeta(s)$ at odd $s \ge 5$ (Theorem~\ref{thm:parity}).
(3)~\textbf{The $\chi_{-4}$ split:} the sub-sums over even and odd $n$
satisfy $D_{\mathrm{even}}(s) - D_{\mathrm{odd}}(s) =
2^{s-1}(\beta(s) - \beta(s{-}2))$ at every integer $s \ge 2$, so Catalan's
$G$ and $\beta(4)$ appear in each half at $s = 4$ and cancel in the sum
(Theorem~\ref{thm:chi4}).
(4)~\textbf{Multi-loop chain sums:} the two-loop sum is $2S_{\min}$, with
$S_{\min}$ in closed form in $\mathbb{Q}[\pi^2, \ln 2, \zeta(3),
\zeta(5)]$ (Hoffman double $t$-values, each reducing to depth $\le 1$);
the three-loop sum $S^{(3)} = 40.6554\ldots$ (rigorous bracket
$[40.65339, 40.65578]$) has an explicit closed form in
$\mathbb{Q}[\pi^2, \ln 2, \zeta(3), \zeta(5), \zeta(5,3)]$, in which
$\zeta(5,3)$, at coefficient $6\pi^2$, is the only constant beyond
single zeta values and every level-2 depth-2 generator cancels; the
four-loop sum has realized depth $\le 3$.  Realized depth is $\le k-1$
at loop order $k = 2, 3, 4$ (Proposition~\ref{prop:depth_k}).
(5)~\textbf{Sommerfeld sums:} the degeneracy-weighted coefficient sums
$D_p$ of the exact Dirac hydrogen spectrum are weight-$(2p{-}1)$ zeta
polynomials in which $\zeta(2)\zeta(2p{-}3)$ cancels and
$\max(0, \lfloor(2p{-}5)/2\rfloor)$ products survive, observed through
$p = 6$ (Observations~\ref{obs:euler_cancel},~\ref{obs:product_rule});
a shell-sum identity removes $\zeta(s{-}1)$ from a fine-structure
Dirichlet series (Theorem~\ref{thm:zeta3_complement}).
The multi-loop sums are defined sums: their vertex weights come from a
vertex model, the scalar-harmonic triangle with a parity clause
(Eq.~\eqref{eq:vertex_rule}).  The physical vector-photon vertex on
$S^3$ obeys a different rule (Paper~33), so the period statements are
statements about the sums as defined.  Two results whose reading depends
on the vertex rule are stated for the model: the ground-state zero of
the model's one-loop self-energy (Theorem~\ref{thm:self_energy_zero}),
which fails for the physical vertex, and an unnormalized one-loop
magnetic-moment ratio (Sec.~\ref{sec:anomalous_moment}), whose
normalization is owed.
\end{abstract}

% ===================================================================
\section{Introduction}
\label{sec:introduction}
% ===================================================================

On the unit three-sphere the Dirac operator has the explicit spectrum
of Camporesi and Higuchi~\cite{camporesi1996}.  Sums over that spectrum
--- spectral zeta values, and the nested sums that arise when
spinor levels are coupled through photon levels in a perturbative
expansion --- are explicit series whose values are periods.  This paper
determines which periods they are.

The questions are those of Paper~18's exchange-constant
taxonomy~\cite{paper18}: which transcendentals enter, through which
structure, and in which ring they sit.  The answers here are exact.
Single spectral sums are governed by operator order
(Secs.~\ref{sec:T9}--\ref{sec:parity}) and by a split of the spectrum
by the parity of $n$ (Sec.~\ref{sec:chi4}).  Nested sums over the
level-2 (alternating) multiple zeta values have closed forms through
three loops and a depth bound at four (Secs.~\ref{sec:smin}--\ref{sec:depth_k}).
Sums over the exact Dirac hydrogen spectrum give Euler sums with a
product-cancellation pattern (Sec.~\ref{sec:euler_cancel}).

\paragraph{Scope.}  The nested sums use a vertex model
(Sec.~\ref{sec:setup}): the triangle rule for scalar harmonics with a
parity clause, weighted by an $SO(4)$ channel count.  They are defined
sums, and every period statement below is a statement about them as
defined.  Whether the vertex model is the vertex of QED on $S^3$ is not
a question this paper answers; the physical vector-photon vertex is
computed in Paper~33, and it is not the model.  Paper~33 also holds the
finite-graph formulations of these diagrams.  The two-term heat trace
and the spectral action built on it are in Paper~51~\cite{paper51}.
All calculations use the \texttt{geovac} modules
\texttt{qed\_vertex}, \texttt{qed\_self\_energy},
\texttt{qed\_anomalous\_moment} and \texttt{qed\_three\_loop}; the
backing tests are named at each result.

% ===================================================================
\section{Spectral Data on $S^3$}
\label{sec:spectra}
% ===================================================================

\subsection{Dirac spectrum}

On the unit three-sphere, the Dirac operator has eigenvalues and
degeneracies~\cite{camporesi1996}
\begin{equation}
  |\lambda_n| = n + \tfrac{3}{2}, \qquad
  g_n = 2(n+1)(n+2), \qquad n = 0, 1, 2, \ldots
  \label{eq:dirac_spectrum}
\end{equation}
The Dirac Dirichlet series is
$D(s) = \sum_{n=0}^{\infty} g_n / |\lambda_n|^s$.
Via Hurwitz zeta,
\begin{equation}
  D(s) = 2\,\zeta(s{-}2, \tfrac{3}{2})
       - \tfrac{1}{2}\,\zeta(s, \tfrac{3}{2}).
  \label{eq:dirac_hurwitz}
\end{equation}
Using $\zeta(s, 3/2) = (2^s - 1)\zeta_R(s) - 2^s$, this becomes
\begin{equation}
  D(s) = 2(2^{s-2} - 1)\,\zeta_R(s{-}2)
       - \tfrac{1}{2}(2^s - 1)\,\zeta_R(s).
  \label{eq:hurwitz_discriminant}
\end{equation}

We define the even-$n$ and odd-$n$ sub-sums:
\begin{equation}
  D_{\mathrm{even}}(s) = \sum_{\substack{n=0 \\ n\,\mathrm{even}}}^{\infty}
    \frac{g_n}{|\lambda_n|^s}, \qquad
  D_{\mathrm{odd}}(s) = \sum_{\substack{n=1 \\ n\,\mathrm{odd}}}^{\infty}
    \frac{g_n}{|\lambda_n|^s},
  \label{eq:D_even_odd_def}
\end{equation}
so that $D(s) = D_{\mathrm{even}}(s) + D_{\mathrm{odd}}(s)$.

\subsection{Photon weights}

The nested sums weight a photon level $q \ge 1$ by
\begin{equation}
  \mu_q = q(q+2), \qquad d_q^T = q(q+2), \qquad q = 1, 2, 3, \ldots
  \label{eq:hodge1_spectrum}
\end{equation}
These are the weights of the sums as defined, and the labelling of
\texttt{geovac/hodge1\_s3.py}.  They are not the Hodge spectrum of the
photon: the coclosed 1-forms on unit $S^3$ at level $q$ have Hodge
eigenvalue $(q+1)^2$, with $q(q+2)$ states in each of the two
$\mathrm{SU}(2)_L \times \mathrm{SU}(2)_R$ components (Ikeda--Taniguchi
1978; Camporesi 1990; recomputed in Paper~33).  The module's anchor test,
\texttt{tests/test\_hodge1\_s3.py}, records the difference.

\subsection{Vertex selection rule}
\label{sec:setup}

The spectral sums of this paper couple spinor harmonics at levels
$n_1, n_2$ to a photon mode at level $q$ under the selection rule
\begin{equation}
  |n_1 - n_2| \le q \le n_1 + n_2, \qquad
  n_1 + n_2 + q \;\text{is odd}.
  \label{eq:vertex_rule}
\end{equation}
This is a \emph{vertex model}, not the vertex of QED on $S^3$: the
triangle bound is the one for scalar harmonics, and the parity clause
is imposed.  The physical vertex
$\int_{S^3} \langle \psi_1, c(A)\,\psi_2 \rangle$, with Clifford
multiplication by a coclosed 1-form $A$, is computed exactly in
Paper~33~\cite{paper33}.  It keeps the model's rule for same-sign spinors inside the
model's window but also opens $q = n_1 + n_2 + 1$, couples spinors of
opposite sign of $D$ at $n_1 + n_2 + q$ even, and couples the ground
state to itself at $q = 1$.  Every sum below that uses
Eq.~\eqref{eq:vertex_rule} is a defined sum, and its period content is
a property of that definition.  The spectral-zeta results of
Theorems~\ref{thm:T9}--\ref{thm:chi4} do not use the rule.

The SO(4) Clebsch--Gordan channel count $W(n_1, n_2, q)$
weights each allowed triple by the number of independent coupling
channels (the two components of the transverse vector harmonic under
$\mathrm{SU}(2)_L \times \mathrm{SU}(2)_R$).  Inside the parity window
the two spinors carry the same content $((n+1)/2, n/2)$, and both
components pass the triangle and integrality conditions in both
factors on every allowed triple, so $W = 2$ there and $0$ elsewhere
(Sec.~\ref{subsec:smin_channel}).

% ===================================================================
\section{Theorem 1: Squared Dirac Spectral Zeta (T9)}
\label{sec:T9}
% ===================================================================

\begin{theorem}[T9: $D^2$ spectral zeta is $\pi^{\mathrm{even}}$]
\label{thm:T9}
The spectral zeta function of the squared Dirac operator on unit
$S^3$ satisfies
\begin{equation}
  \zeta_{D^2}(s)
    = 2^{2s-1}\bigl[\lambda(2s{-}2) - \lambda(2s)\bigr],
  \label{eq:T9}
\end{equation}
where $\lambda(2k) = (1 - 2^{-2k})\,\zeta_R(2k)
= \mathrm{rational} \times \pi^{2k}$.
At every integer~$s \ge 1$, $\zeta_{D^2}(s)$ is a two-term polynomial
in $\pi^2$ with rational coefficients;\ no odd-zeta content appears
at any integer $s \ge 1$.  The statement is an integer-$s$ one:\ at
half-integer $s \ge 5/2$ the same formula~\eqref{eq:T9} carries odd zeta
values (e.g.\ $\zeta_{D^2}(5/2) = 14\,\zeta(3) - \tfrac{31}{2}\,\zeta(5)$),
and $s = 1/2, 3/2$ are simple poles with rational residues.
\end{theorem}

\begin{proof}
The squared Dirac eigenvalues are $\lambda_n^2 = (n + 3/2)^2
= (2n+3)^2/4$, so
$\zeta_{D^2}(s) = \sum_{n \ge 0} g_n \cdot (4/(2n+3)^2)^s
= 4^s \sum_{n \ge 0} 2(n+1)(n+2) / (2n+3)^{2s}$.
Substituting $m = 2n+3$ (which runs over odd integers $\ge 3$), the
sum becomes $\sum_{m \ge 3, m\,\mathrm{odd}} (m^2-1)/2 \cdot
4^s / m^{2s}$.  The $(m^2 - 1)/2$ factor splits as
$(m^{2-2s} - m^{-2s})/2$, giving two Dirichlet-type sums over odd
integers.  Each odd-integer sum is $\lambda(2k) = (1-2^{-2k})\zeta_R(2k)$
with the first three odd values ($m=1$) added or subtracted to correct
the lower limit.  By the Euler--Bernoulli formula,
$\zeta_R(2k) = (-1)^{k+1} (2\pi)^{2k} B_{2k} / (2\cdot (2k)!)$,
so $\lambda(2k) = \mathrm{rational} \times \pi^{2k}$.  At integer $s$
both arguments $2s-2$ and $2s$ are even, so odd-zeta values
$\zeta_R(2k+1)$ cannot appear;\ at half-integer $s \ge 5/2$ both
arguments are odd and $\ge 3$, so $\lambda(2s-2)$, $\lambda(2s)$ are
rational multiples of odd zeta values (at $s = 1/2, 3/2$ an argument
equals $1$, a pole;\ at negative half-integer $s$ the values
$\zeta(1-2k) = -B_{2k}/(2k)$ are rational).
\end{proof}

\noindent\textbf{Verification:} Symbolic identity confirmed in
\texttt{sympy} for $s = 1, \ldots, 10$.  Numerical agreement to
$10^{-60}$ at $s = 2, \ldots, 20$.  The integer-$s$ closed forms at
$s = 2, 3$ and the half-integer value
$\zeta_{D^2}(5/2) = 14\,\zeta(3) - \tfrac{31}{2}\,\zeta(5)$ are checked
against the direct spectral sum in
\texttt{tests/test\_paper28\_t9\_half\_integer.py}.

\begin{table}[h]
\centering
\caption{Squared Dirac spectral zeta $\zeta_{D^2}(s)$.}
\label{tab:D2_values}
\begin{tabular}{@{}cl@{}}
\toprule
$s$ & $\zeta_{D^2}(s)$ \\
\midrule
1 & $-\pi^2/4$ \\
2 & $\pi^2 - \pi^4/12$ \\
3 & $\pi^4/3 - \pi^6/30$ \\
4 & $2\pi^6/15 - 17\pi^8/1260$ \\
\bottomrule
\end{tabular}
\end{table}

\begin{remark}[Where $\sqrt{\pi}$ lives:\ heat-trace coefficients,
not $\zeta_{D^2}$]
The $\sqrt{\pi}$ met in heat-kernel-side analyses of $\zeta_{D^2}$
does not live in the spectral zeta function at any $s$.  It lives only
in the heat-trace coefficients:\ the short-time expansion
$\mathrm{Tr}\,e^{-tD^2} \sim \tfrac{\sqrt{\pi}}{2}\,t^{-3/2}
- \tfrac{\sqrt{\pi}}{4}\,t^{-1/2}$ carries the $(4\pi t)^{-3/2}$-type
prefactor of an odd-dimensional heat kernel, and the Mellin transform
$\zeta_{D^2}(s) = \Gamma(s)^{-1}\int_0^\infty t^{s-1}\,
\mathrm{Tr}\,e^{-tD^2}\,dt$ divides each coefficient by $\Gamma(s)$ at
its pole.  At $s = 3/2$ this gives the residue
$(\sqrt{\pi}/2)/\Gamma(3/2) = 1$, and at $s = 1/2$ the residue
$(-\sqrt{\pi}/4)/\Gamma(1/2) = -1/4$:\ both rational.  Away from the
poles, the values at half-integer $s \ge 5/2$ are odd-zeta combinations
(Theorem~\ref{thm:T9}), not multiples of $\sqrt{\pi}$.  At integer
$s \ge 1$ the values are pure $\pi^{2k}\cdot\mathbb{Q}$
(Table~\ref{tab:D2_values} above). After
the standard Connes--Chamseddine volume normalization (multiplication
of the heat-trace coefficients by $(4\pi)^{d/2}$ with $d=3$), the
coefficients also cleanly enter $\mathbb{Q}[\pi^2]$:\ Paper~51~\cite{paper51} Cor.~2.1
gives $a_0^{D^2} = 4\pi^2$, $a_1^{D^2} = -2\pi^2$, $a_k = 0$ for
$k \ge 2$. At integer $s$ the GeoVac $S^3$ M2 sub-mechanism is therefore mixed-Tate
over $\mathbb{Q}$ in the F--M sense (Paper~32~\cite{paper32}
\S~VIII Corollary), in the strictly smaller pure-Tate sub-ring
$\bigoplus_k \pi^{2k}\cdot\mathbb{Q}$.
\end{remark}

% ===================================================================
\section{Theorem 2: Even/Odd Parity Discriminant}
\label{sec:parity}
% ===================================================================

\begin{theorem}[Parity discriminant]
\label{thm:parity}
At even integer $s$, the Dirac Dirichlet series $D(s)$ is purely
$\pi^{\mathrm{even}}$.  At odd integer $s$, $D(s)$ involves
$\zeta_R(s{-}2)$ and $\zeta_R(s)$ (odd-zeta values;\ $s \ge 5$, since
$\zeta_R(1)$ is a pole).  Operator order (2nd vs 1st) is the primary
transcendental discriminant, conditional on the period conjecture
$\zeta(2k{+}1) \notin \mathbb{Q}(\pi)$:\ the second-order operators
$\Delta_{\mathrm{LB}}$ and $D^2$ produce
$\pi^{\mathrm{even}}$ exclusively at integer $s$ (this holds for $\Delta_{\mathrm{LB}}$ and $D^2$, with spectra $n^2 - 1$
and $(n+3/2)^2$, whose partial fractions have rational poles;\ it is not a property of every second-order operator ---
the spinor connection Laplacian $\nabla^*\nabla = D^2 - 3/2$ has
$\zeta_{\nabla^*\nabla}(2) = \tfrac{1}{12}\bigl(14\,\pi\tan(\tfrac{\pi\sqrt6}{2})/\sqrt6
+ 5\pi^2\sec^2(\tfrac{\pi\sqrt6}{2})\bigr)$, outside $\mathbb{Q}[\pi^2]$) (their spectral
determinants are not covered:\ $\zeta'_{D^2}(0) = 3\zeta_R(3)/(4\pi^2) + \tfrac12\ln 2$); first-order
operators produce odd-zeta values at odd~$s$.
\end{theorem}

\begin{proof}
From Eq.~\eqref{eq:hurwitz_discriminant},
$D(s) = 2(2^{s-2} - 1)\,\zeta_R(s{-}2)
       - \tfrac{1}{2}(2^s - 1)\,\zeta_R(s)$.
At even~$s$, both $\zeta_R(s{-}2)$ and $\zeta_R(s)$ are Bernoulli-type:
$\zeta_R(2k) = (-1)^{k+1} (2\pi)^{2k} B_{2k} / (2 \cdot (2k)!)
= \mathrm{rational} \times \pi^{2k}$.  At odd~$s$, both
$\zeta_R(s{-}2)$ and $\zeta_R(s)$ are genuine odd-zeta values
($\zeta(3)$, $\zeta(5)$, \ldots), conjecturally $\mathbb{Q}$-linearly
independent of all $\pi^{2k}$ (open:\ Ap\'ery's theorem gives the
irrationality of $\zeta(3)$ and Rivoal's infinitely many irrational odd
zeta values~\cite{apery1979,rivoal2000};\ neither gives independence
from $\pi$).
\end{proof}

\begin{table}[h]
\centering
\caption{Dirac Dirichlet series $D(s)$: PSLQ identifications at 80 dps.}
\label{tab:D_values}
\begin{tabular}{@{}cll@{}}
\toprule
$s$ & $D(s)$ & Class \\
\midrule
4 & $\pi^2 - \pi^4/12$ & $\pi^{\mathrm{even}}$ \\
5 & $14\,\zeta(3) - \tfrac{31}{2}\,\zeta(5)$ & odd-zeta \\
6 & $\pi^4/3 - \pi^6/30$ & $\pi^{\mathrm{even}}$ \\
7 & $62\,\zeta(5) - \tfrac{127}{2}\,\zeta(7)$ & odd-zeta \\
8 & $2\pi^6/15 - 17\pi^8/1260$ & $\pi^{\mathrm{even}}$ \\
9 & $254\,\zeta(7) - \tfrac{511}{2}\,\zeta(9)$ & odd-zeta \\
10 & $17\pi^8/315 - 31\pi^{10}/5670$ & $\pi^{\mathrm{even}}$ \\
\bottomrule
\end{tabular}
\end{table}

% ===================================================================
\section{Theorem 3: The $\chi_{-4}$ Identity}
\label{sec:chi4}
% ===================================================================

\begin{theorem}[$\chi_{-4}$ identity]
\label{thm:chi4}
For every integer $s \ge 2$ (the even/odd spectral sums converge for $s \ge 4$;\ at $s = 2, 3$ by Hurwitz continuation),
\begin{equation}
  D_{\mathrm{even}}(s) - D_{\mathrm{odd}}(s)
    \;=\; 2^{s-1}\bigl(\beta(s) - \beta(s{-}2)\bigr),
  \label{eq:D_diff_closed}
\end{equation}
where $\beta(s) = L(s, \chi_{-4})$ is the Dirichlet beta function.
\end{theorem}

\begin{proof}
Express $D_{\mathrm{even}}(s)$ and $D_{\mathrm{odd}}(s)$ via Hurwitz
zeta at quarter-integer shifts.  With $m = n + 3/2$ and full
degeneracy $g_n = 2(n+1)(n+2) = 2(m^2 - 1/4)$, the even-$n$
eigenvalues are $m = 2(k + 3/4)$ and the odd-$n$ ones
$m = 2(k + 5/4)$, $k \ge 0$, so
$D_{\mathrm{even}}(s) = 2^{-(s-3)} \zeta(s{-}2, 3/4)
                       - 2^{-(s+1)} \zeta(s, 3/4)$
and
$D_{\mathrm{odd}}(s) = 2^{-(s-3)} \zeta(s{-}2, 5/4)
                       - 2^{-(s+1)} \zeta(s, 5/4)$.
Two elementary Hurwitz identities:
$\zeta(s, 5/4) = \zeta(s, 1/4) - 4^s$ (shift recurrence) and
$\zeta(s, 3/4) - \zeta(s, 1/4) = -4^s \beta(s)$ (the Dirichlet
beta definition).  Together,
$\zeta(s, 3/4) - \zeta(s, 5/4) = 4^s(1 - \beta(s))$;\ substituting
into the difference,
$D_{\mathrm{even}} - D_{\mathrm{odd}}
 = 2^{s-1}(1 - \beta(s{-}2)) - 2^{s-1}(1 - \beta(s))
 = 2^{s-1}(\beta(s) - \beta(s{-}2))$,
which is~\eqref{eq:D_diff_closed}.  The Hurwitz decompositions and
the closed form are verified bit-exactly against direct spectral sums
at $s \in \{4, 5, 6, 7\}$ to $40$ digits.
\end{proof}

\noindent\textbf{Verification:}
Eq.~\eqref{eq:D_diff_closed} verified symbolically (\texttt{sympy}
residual zero) at $s \in \{2, 3, \ldots, 10, 15\}$ and by PSLQ at
100-digit precision at every integer $s$ in the tested range.  At $s=4$,
$D_{\mathrm{even}}(4) - D_{\mathrm{odd}}(4) = 8(\beta(4) - G)$
(since $G = \beta(2)$), consistent with the explicit
decompositions~\eqref{eq:D_even_4}--\eqref{eq:D_odd_4} below.

\subsection{The even/odd split at $s = 4$}

The even-$n$ and odd-$n$ sub-sums of $D(4)$ decompose as:
\begin{align}
  D_{\mathrm{even}}(4)
    &= \tfrac{\pi^2}{2} - \tfrac{\pi^4}{24}
       - 4G + 4\beta(4),
  \label{eq:D_even_4} \\
  D_{\mathrm{odd}}(4)
    &= \tfrac{\pi^2}{2} - \tfrac{\pi^4}{24}
       + 4G - 4\beta(4),
  \label{eq:D_odd_4} \\
  D(4) &= \pi^2 - \tfrac{\pi^4}{12}
  \qquad\text{(Catalan and $\beta(4)$ cancel)}.
  \label{eq:D_full}
\end{align}
Here $G = \beta(2) = \sum_{n=0}^{\infty} (-1)^n/(2n+1)^2 \approx
0.9159656$ is Catalan's constant and $\beta(4) \approx 0.9889446$ is the
Dirichlet beta function at $s=4$.  The mechanism is the
quarter-integer Hurwitz shift: restricting $n$ to one parity class puts
the eigenvalues $m = n + 3/2$ on $2(k + 3/4)$ or $2(k + 5/4)$, and
$L(s, \chi_{-4})$ appears in the difference.  In the vertex model of
Sec.~\ref{sec:setup} a parity restriction of this kind is what the
parity clause imposes on internal lines; that reading belongs to the
model.  Verified by PSLQ at 80-digit precision with basis
$\{1, \pi^2, \pi^4, G, \beta(4)\}$ (\texttt{tests/test\_smin\_chi\_neg4.py},
\texttt{tests/test\_qed\_vertex.py}).

\begin{remark}[Motivic home of Theorem~\ref{thm:chi4}]
\label{rem:chi4_motivic}
Each of $D_{\mathrm{even}}(s)$, $D_{\mathrm{odd}}(s)$, and their
difference is a period of mixed Tate motives at level $4$ over
$\mathcal{O}_4[1/4] = \mathbb{Z}[i, 1/2]$, by direct application of
Deligne's classification at $N = 4$ and the explicit basis of
Glanois (Paper~32~\S~VIII Corollary~\texttt{cor:m3\_cyclotomic\_mixed\_tate};
Paper~18~\S~III.7; \texttt{tests/test\_paper55\_m3\_cyclotomic\_mixed\_tate.py}).
The split can be read as a Galois descent from level $4$ to level $2$:\
the un-restricted
$D(s) = D_{\mathrm{even}} + D_{\mathrm{odd}}$ lives in the smaller
level-$\le 2$ ring (Theorem~\ref{thm:T9} and
Theorem~\ref{thm:parity}); the $\zeta(s, 1/4)$ pieces in
$D_{\mathrm{even}}, D_{\mathrm{odd}}$ individually cancel in the
sum but survive in the difference, exposing the level-$4$ Dirichlet
$L$-value content.  Catalan $G = \beta(2)$ and $\beta(4)$ are
periods in $\mathcal{MT}(\mathbb{Z}[i, 1/2])$ at depth $1$ but
conjecturally not in $\mathcal{MT}(\mathbb{Z}[1/2])$ at any
depth, so the level-$4$ content is carried by the parity split.
\end{remark}

\subsection{Three axes}

The single spectral sums of Secs.~\ref{sec:T9}--\ref{sec:chi4} are
sorted by three independent axes:
\begin{enumerate}
\item \textbf{Operator order.}
      Second-order ($D^2$) $\to$ $\pi^{\mathrm{even}}$ at integer $s$ (Theorem~\ref{thm:T9});
      first-order ($|D|$) $\to$ odd-zeta at odd $s \ge 5$ (Theorem~\ref{thm:parity}).
\item \textbf{Bundle type.}
      Scalar vs spinor modulates coefficients but not the
      transcendental class (Paper~18 Sec~IV).
\item \textbf{Restriction of the spectrum.}
      Restricting $n$ to one parity class introduces the Dirichlet
      $L$-values $G$, $\beta(4)$, absent from the unrestricted sums
      (Theorem~\ref{thm:chi4}).
\end{enumerate}

% ===================================================================
\section{Theorem 4: Ground-State Zero of the Model Self-Energy}
\label{sec:self_energy}
% ===================================================================

The model's one-loop self-energy sum is
\begin{equation}
  \Sigma(n_{\mathrm{ext}}) = \sum_{n_{\mathrm{int}}, q}
    \frac{W(n_{\mathrm{ext}}, n_{\mathrm{int}}, q) \;
          g_{n_{\mathrm{int}}} \; d_q^T}
         {|\lambda_{n_{\mathrm{int}}}|^{2s_e} \; \mu_q^{s_\gamma}},
  \label{eq:self_energy}
\end{equation}
where the sum is over internal electron mode $n_{\mathrm{int}}$ and
photon mode $q$ subject to the vertex selection
rule~\eqref{eq:vertex_rule}, $s_e = 2$ (so $|\lambda|^4$ in the
denominator), $s_\gamma = 1$, and $W$ is the SO(4) channel count.

\begin{theorem}[Self-energy zero of the vertex model]
\label{thm:self_energy_zero}
Under the selection rule~\eqref{eq:vertex_rule}, $\Sigma(0) = 0$
exactly: the model's one-loop self-energy vanishes identically at the
Dirac ground state $n_{\mathrm{ext}} = 0$.
\end{theorem}

\begin{proof}
At $n_{\mathrm{ext}} = 0$, the vertex selection rule requires
$|0 - n_{\mathrm{int}}| \le q \le 0 + n_{\mathrm{int}}$ and
$0 + n_{\mathrm{int}} + q$ odd.  The first condition forces
$q = n_{\mathrm{int}}$.  The second then requires
$2\,n_{\mathrm{int}}$ to be odd, which is impossible for any
integer $n_{\mathrm{int}}$.  Therefore $W(0, n_{\mathrm{int}}, q) = 0$
for all $(n_{\mathrm{int}}, q)$, and the sum is empty.
\end{proof}

\noindent The theorem is a statement about the model; it is backed by
\texttt{tests/test\_qed\_self\_energy.py}.  The zero comes from the
model's triangle bound, which is the scalar one.  It fails for the
physical vector-photon vertex: there the ground state couples at
$q = n_{\mathrm{int}} + 1$ to same-sign internal levels and at
$q = n_{\mathrm{int}}$ to opposite-sign ones, so the corresponding sum
is a sum of positive terms with terms present (Paper~33).

\subsection{Self-energy for excited states}

For $n_{\mathrm{ext}} \ge 1$, the model self-energy is nonzero and
convergent at $s_e = 2$, $s_\gamma = 1$, and increases monotonically
with $n_{\mathrm{ext}}$.

\subsection{Vertex correction}

The one-loop vertex correction involves two vertex insertions:
\begin{equation}
  \Lambda(n_{\mathrm{ext}}) = \sum_{n_{\mathrm{int}}, q_1, q_2}
    \frac{W(n_{\mathrm{ext}}, n_{\mathrm{int}}, q_1) \;
          W(n_{\mathrm{int}}, n_{\mathrm{ext}}, q_2) \;
          g_{n_{\mathrm{int}}} \; d_{q_1}^T \; d_{q_2}^T}
         {|\lambda_{n_{\mathrm{int}}}|^{2s_e} \;
          \mu_{q_1}^{s_\gamma} \; \mu_{q_2}^{s_\gamma}}.
  \label{eq:vertex_correction}
\end{equation}
This is $O(N^3)$ in the cutoff.  Both $\Lambda(n_{\mathrm{ext}})$ and
$\Sigma(n_{\mathrm{ext}})$ are backed by 35 tests in
\texttt{tests/test\_qed\_self\_energy.py}.

\noindent\textbf{Note:} The shell-summed vertex correction
$\Lambda(n_{\mathrm{ext}})$ defined above uses the CG-weight channel
count $W$ and is $m_j$-independent (it gives $B = 0$ by the
Wigner--Eckart theorem).  The anomalous magnetic moment extraction
requires the $m_j$-resolved three-point function developed in the
next section.

% ===================================================================
\section{A One-Loop Magnetic-Moment Ratio in the Vertex Model}
\label{sec:anomalous_moment}
% ===================================================================

This section is a statement about the vertex model.  Its sums carry
unit coupling, no loop measure and the photon weights of
Eq.~\eqref{eq:hodge1_spectrum}, and the loop vertex obeys
Eq.~\eqref{eq:vertex_rule}.  Normalizing them --- the coupling, the
transverse photon propagator with its mode measure and the Hodge
eigenvalue $(q+1)^2$, and the physical vertex of Paper~33 --- is owed
work, so no number below is an anomalous moment.  The algebraic
content (Clebsch--Gordan sums over rational denominators) is what the
section records.

\subsection{Why the two-point vertex correction gives $B = 0$}

The two-point vertex correction $\Lambda(n_{\mathrm{ext}})$ defined
in Eq.~\eqref{eq:vertex_correction} sums over all internal modes and
photon polarizations at \emph{fixed} external quantum numbers
$(n_{\mathrm{ext}}, j_{\mathrm{ext}}, m_j)$.  On $S^3$, the vertex
amplitude at each junction involves a product of four
$\mathrm{SU}(2)_L \times \mathrm{SU}(2)_R$ Clebsch--Gordan
coefficients.  When the photon polarizations $(m_{g,L}, m_{g,R})$ are
summed over at both vertices, the SO(4) Wigner--Eckart theorem forces
the result to be independent of~$m_j$:
\begin{equation}
  \Lambda(n_{\mathrm{ext}}, j, m_j = +\tfrac{1}{2})
  = \Lambda(n_{\mathrm{ext}}, j, m_j = -\tfrac{1}{2}).
  \label{eq:B_zero_two_point}
\end{equation}
Therefore the spin-dependent piece
$B_{\mathrm{2pt}} = \Lambda(+\tfrac{1}{2}) - \Lambda(-\tfrac{1}{2}) = 0$
identically.  The two-point vertex correction is a scalar under SO(4)
rotations and cannot distinguish spin orientations.  To extract the
anomalous magnetic moment ($F_2$ form factor), an external probe must
break the rotational symmetry.

\subsection{Three-point vertex correction with probe photon}

The anomalous magnetic moment is extracted from a three-point function
in which an external probe photon of mode $q_{\mathrm{probe}} = 1$
(the lowest transverse vector harmonic, the $S^3$ analog of a uniform
magnetic field) is inserted between the two loop-photon vertices.  The
topology is:
\begin{equation}
  \langle \mathrm{ext}, m_j |
  \;\Gamma^{\mathrm{1-loop}}_{\mathrm{probe}}\;
  | \mathrm{ext}, m_j \rangle
  = \sum_{\substack{n_{\mathrm{int}},\, j_{\mathrm{int}} \\
                     m_{j,\mathrm{int}},\, m'_{j,\mathrm{int}} \\
                     q_{\mathrm{loop}}}}
    \frac{V_1 \;\cdot\; V_{\mathrm{probe}} \;\cdot\; V_3}
         {|\lambda_{n_{\mathrm{int}}}|^4 \;\mu_{q_{\mathrm{loop}}}},
  \label{eq:three_point_vertex}
\end{equation}
where:
\begin{itemize}
\item $V_1$: amplitude for the external electron $(m_j)$ to emit a
      loop photon and transition to the internal state $(m_{j,\mathrm{int}})$;
\item $V_{\mathrm{probe}}$: amplitude for the probe photon at $q=1$ to
      scatter the internal electron from $m_{j,\mathrm{int}}$ to
      $m'_{j,\mathrm{int}}$, summed over probe polarizations;
\item $V_3$: amplitude for the internal state $(m'_{j,\mathrm{int}})$ to
      reabsorb the loop photon and return to the external state $(m_j)$.
\end{itemize}
Each vertex amplitude is a product of four Clebsch--Gordan coefficients
coupling the $\mathrm{SU}(2)_L \times \mathrm{SU}(2)_R$ representations
of source, target, and photon:
\begin{equation}
  V = \sum_{m_{L,1}}
    C^{j_s}_{j_{s,L}\, j_{s,R}}(m_{L,1},\, m_j - m_{L,1};\, m_j)
    \;\cdot\;
    C^{j_t}_{j_{t,L}\, j_{t,R}}(m_{L,2},\, m_{R,2};\, m'_j)
    \;\cdot\;
    C^{j_{t,L}}_{j_{s,L}\, j_{g,L}}(m_{L,1},\, m_{g,L};\, m_{L,2})
    \;\cdot\;
    C^{j_{t,R}}_{j_{s,R}\, j_{g,R}}(m_{R,1},\, m_{g,R};\, m_{R,2}),
  \label{eq:4CG_vertex}
\end{equation}
where $m_{L,2} = m_{L,1} + m_{g,L}$ and $m_{R,2} = m_{R,1} + m_{g,R}$,
and the sum runs over all $m_{L,1} \in \{-j_{s,L}, \ldots, +j_{s,L}\}$.

The probe photon at $q = 1$ has channels $(j_{g,L}, j_{g,R}) = (1, 0)$
or $(0, 1)$, corresponding to the two components of the transverse
vector harmonic.  The probe couples internal states
$m_{j,\mathrm{int}} \to m'_{j,\mathrm{int}}$, enabling spin-dependent
amplitudes that the polarization-summed two-point function cannot access.

The spin-dependent piece of the vertex correction is:
\begin{equation}
  B = \Gamma^{\mathrm{1-loop}}_{\mathrm{probe}}(m_j = +\tfrac{1}{2})
    - \Gamma^{\mathrm{1-loop}}_{\mathrm{probe}}(m_j = -\tfrac{1}{2}).
  \label{eq:B_three_point}
\end{equation}

\subsection{Tree-level normalization}

The tree-level (zeroth-order) coupling of the external electron to the
$q = 1$ probe photon provides the normalization for $F_2$ extraction.
At the external state $n_{\mathrm{ext}} = 1$, $j = 1/2$
(Camporesi--Higuchi level $n_{\mathrm{CH}} = 1$, with
$j_L = 1$, $j_R = 1/2$):
\begin{align}
  V_{\mathrm{tree}}(m_j = +\tfrac{1}{2})
    &= -V_{\mathrm{tree}}(m_j = -\tfrac{1}{2})
    = 0.2790,
  \label{eq:V_tree} \\
  V_{\mathrm{magnetic}}
    &= V_{\mathrm{tree}}(+\tfrac{1}{2})
     - V_{\mathrm{tree}}(-\tfrac{1}{2})
    = 0.5579.
  \label{eq:V_magnetic}
\end{align}
The charge piece
$[V_{\mathrm{tree}}(+\tfrac{1}{2}) + V_{\mathrm{tree}}(-\tfrac{1}{2})]/2
= 0$
vanishes---the $q = 1$ probe photon couples purely to the magnetic
moment at tree level, as expected for the $S^3$ analog of a uniform
magnetic field.

Analytically, $V_{\mathrm{magnetic}} = 2(\sqrt{2}/3 - \sqrt{3}/9)$,
a pure algebraic number from the Clebsch--Gordan coefficients.

The Pauli form factor is then:
\begin{equation}
  F_2 = \frac{B}{V_{\mathrm{magnetic}}},
  \label{eq:F2_def}
\end{equation}
and the anomalous magnetic moment on $S^3$ is $a_e^{S^3} = F_2$.

\subsection{Convergence and results}

The three-point vertex correction converges rapidly.  Table~\ref{tab:g_minus_2}
shows the per-level contributions for the external state at
$n_{\mathrm{ext}} = 1$.

\begin{table}[h]
\centering
\caption{Per-level contributions to the anomalous magnetic moment on
$S^3$.  External state $n_{\mathrm{ext}} = 1$, $j = 1/2$.
The $n_{\mathrm{int}} = 0$ contribution vanishes in the vertex model
(Theorem~\ref{thm:self_energy_zero}); with the vector-photon rule the
$q = 1$ channel there is open.
Contributions drop by approximately $100\times$ per successive allowed
level pair.}
\label{tab:g_minus_2}
\begin{tabular}{@{}ccccc@{}}
\toprule
$n_{\mathrm{int}}$ & $\Delta B$ & $B_{\mathrm{cumul}}$
  & $F_2 / [\alpha/(2\pi)]$ & Status \\
\midrule
0 & $0$     & $0$           & ---     & closed in the vertex model \\
1 & $6.96 \times 10^{-4}$  & $6.96 \times 10^{-4}$  & 1.074  & 99\% of final \\
2 & $3.76 \times 10^{-6}$  & $7.00 \times 10^{-4}$  & 1.080  & \\
3 & $2.85 \times 10^{-6}$  & $7.02 \times 10^{-4}$  & 1.084  & \\
4 & $1.34 \times 10^{-7}$  & $7.025 \times 10^{-4}$ & 1.084  & \\
5 & $1.25 \times 10^{-7}$  & $7.027 \times 10^{-4}$ & 1.084  & \\
6 & $1.25 \times 10^{-8}$  & $7.027 \times 10^{-4}$ & 1.084  & \\
7 & $1.36 \times 10^{-8}$  & $7.027 \times 10^{-4}$ & 1.084  & \\
8 & $2.00 \times 10^{-9}$  & $7.027 \times 10^{-4}$ & 1.084  & \\
\bottomrule
\end{tabular}
\end{table}

Three structural features are visible:
\begin{enumerate}
\item \textbf{Rapid convergence.}  The $n_{\mathrm{int}} = 1$ level
      contributes 99\% of the final value.  Each subsequent allowed level
      pair contributes roughly $100\times$ less than the previous one.
\item \textbf{Vertex parity filtering (model).}
      The $n_{\mathrm{int}} = 0$ contribution vanishes by the model
      rule~\eqref{eq:vertex_rule} (the loop-photon vertex at
      $n_{\mathrm{ext}} = 1$, $n_{\mathrm{int}} = 0$ requires
      $q_{\mathrm{loop}} = 1$ and $1 + 0 + 1 = 2$ is even, forbidden).
      The dominant contribution comes from $n_{\mathrm{int}} = 1$
      (the nearest allowed level to the external state).
      Higher even-$n_{\mathrm{int}}$ levels contribute but are
      suppressed relative to odd-$n_{\mathrm{int}}$ neighbors.
\item \textbf{The ratio.}  The ratio $B/V_{\mathrm{mag}}$ divided by
      the CODATA value of $\alpha/(2\pi)$ converges to $1.084$.
\end{enumerate}

\subsection{The ratio at $n_{\mathrm{ext}} = 1$}

At $n_{\max} = 12$ (converged to 5 significant figures; the value is
already stable at $n_{\max} = 8$):
\begin{equation}
  \frac{F_2}{\alpha/(2\pi)} = 1.084,
  \label{eq:F2_ratio}
\end{equation}
with $B = 7.027 \times 10^{-4}$ and
$V_{\mathrm{magnetic}} = 0.5579$.

\noindent\textbf{Normalization [OBSERVATION].}  $B$ and
$V_{\mathrm{mag}}$ are sums of Clebsch--Gordan products over spectral
denominators with unit coupling: neither carries $e^2$, a loop measure
or the harmonic normalization on $S^3$.  The ratio
$F_2 = B/V_{\mathrm{mag}} = 1.2596 \times 10^{-3}$ is therefore not an
anomalous moment in units of $\alpha$, and Eq.~\eqref{eq:F2_ratio}
divides it by the CODATA value of $\alpha/(2\pi)$; the coupling does
not cancel.  The value $1.084$ is a numerical observation about this
unnormalized ratio, computed moreover in the vertex model (the
$n_{\mathrm{int}} = 0$, $q = 1$ channel that the physical vertex opens
is closed).  The flat-space limit $R \to \infty$ is not computed.

\subsection{Mode dependence}
\label{sec:mode_dependence}

Table~\ref{tab:mode_dependence} shows the same ratio at
$n_{\mathrm{ext}} = 1, \ldots, 5$.

\begin{table}[h]
\centering
\caption{The model's magnetic-moment ratio on unit $S^3$ by external
level.  Each row reports the tree-level magnetic coupling
$V_{\mathrm{mag}}$, the one-loop correction $B$, and $B/V_{\mathrm{mag}}$
divided by the CODATA $\alpha/(2\pi)$ (unnormalized; see above).
The $n_{\mathrm{ext}} = 1$ result uses $n_{\max} = 12$ (converged
to 5 digits); $n_{\mathrm{ext}} = 2$--$5$ use $n_{\max} = 5$
(converged to 3 digits via $n^{-7}$ power-law tail).
Exact algebraic forms for $V_{\mathrm{mag}}$ and the dominant
$B(n_{\mathrm{int}}=1)$ are available at each $n_{\mathrm{ext}}$;
all expressions are sums of square roots with rational coefficients
(no transcendentals).}
\label{tab:mode_dependence}
\begin{tabular}{@{}cccccr@{}}
\toprule
$n_{\mathrm{ext}}$ & $\lambda_{\mathrm{ext}}$
  & $V_{\mathrm{mag}}$ & $B$
  & ratio to CODATA $\alpha/(2\pi)$ & ratio $-1$ \\
\midrule
1 & $5/2$   & $0.5579$ & $7.027 \times 10^{-4}$ & $1.084$ & $+8.4\%$ \\
2 & $7/2$   & $0.3893$ & $2.772 \times 10^{-4}$ & $0.613$ & $-39\%$ \\
3 & $9/2$   & $0.3001$ & $1.296 \times 10^{-4}$ & $0.370$ & $-63\%$ \\
4 & $11/2$  & $0.2445$ & $6.940 \times 10^{-5}$ & $0.244$ & $-76\%$ \\
5 & $13/2$  & $0.2064$ & $4.139 \times 10^{-5}$ & $0.173$ & $-83\%$ \\
\bottomrule
\end{tabular}
\end{table}

The loop correction $B$ drops faster with $n_{\mathrm{ext}}$ than the
tree-level coupling: across the five levels $B$ spans a factor of
$17\times$ while $V_{\mathrm{mag}}$ spans $2.7\times$.  The tree-level
amplitude follows $V_{\mathrm{mag}} \to (4/3)/|\lambda|$
asymptotically (Richardson-extrapolated to 4-digit agreement with
$4/3$).  A leading curvature form $1 + R_{\mathrm{scalar}}/(12|\lambda|^2)$
(motivated by the curved-space formalism of~\cite{parker1980}) gives
$1.080$ at $n_{\mathrm{ext}} = 1$; given the normalization above, the
closeness to $1.084$ is a numerical coincidence, and a two-term fit in
$1/|\lambda|^2$ to the five levels gives a leading coefficient near
$-20$, not $1/2$.  With one curvature value per level, coefficients
beyond the first cannot be separated.

The five computed values are well described by a power law:
\begin{equation}
  \frac{F_2}{\alpha/(2\pi)} = (5.91 \pm 0.36)\;|\lambda|^{-1.84 \pm 0.06}.
  \label{eq:mode_power_law}
\end{equation}
The fitted exponent is $2.8\sigma$ from $-2$; the forced $p = -2$ fit
has twice the RMS residual.  Analysis of the dominant internal-mode
contribution ($n_{\mathrm{int}} = 1$) shows its local log-log slope
converging toward $-2$ ($-2.57 \to -2.26 \to -2.16 \to -2.11$ across
$n_{\mathrm{ext}} = 1, \ldots, 5$), but higher internal modes
contribute an increasing fraction at larger $n_{\mathrm{ext}}$
($99\%$ at $n_{\mathrm{ext}} = 1$ vs.\ $69\%$ at $n_{\mathrm{ext}} = 3$),
pulling the effective exponent above $-2$.  Whether the true asymptotic
exponent is exactly $-2$ (with finite-curvature corrections) or a
distinct value remains an open question.

The flat-space limit is $R \to \infty$ at fixed physical mass
($n \sim m R \to \infty$ simultaneously), not $|\lambda| \to \infty$ on
a fixed-size sphere, and it is not computed here.

\medskip
\noindent\textbf{Verification:} The computation uses exact
\texttt{sympy} Clebsch--Gordan coefficients at every vertex, with
no numerical approximation in the angular algebra.  The propagator
denominators are exact rationals ($|\lambda_n|^4$ and $\mu_q$ from
Eqs.~\eqref{eq:dirac_spectrum} and~\eqref{eq:hodge1_spectrum}).
Backing tests: \texttt{tests/test\_qed\_anomalous\_moment.py}.

% ===================================================================
\section{Heat Trace and Spectral Action}
\label{sec:one_loop}
% ===================================================================

The Dirac heat trace on unit $S^3$ is two-term up to exponentially
small corrections,
\begin{equation}
  \operatorname{Tr} e^{-tD^2}
  = \frac{\sqrt{\pi}}{2}\,t^{-3/2} - \frac{\sqrt{\pi}}{4}\,t^{-1/2}
  + O(e^{-\pi^2/t}),
  \label{eq:sd_two_term}
\end{equation}
a result of Chamseddine and Connes~\cite{chamseddine_connes_uncanny}.
Equivalently $\zeta_{D^2}(-k) = 0$ for every integer $k \ge 0$, the
negative-integer companion of Theorem~\ref{thm:T9}.  The proof by
Bernoulli polynomials at $3/2$, the spectral action on $S^3_R$ and on
$S^3 \times S^1$ that follows from it, the scalar and
transverse-traceless comparison heat traces, and the graviton
diagnostics on the truncated spectrum are in Paper~51~\cite{paper51}.
The cutoff $\Lambda$ defined by $\operatorname{Tr}
e^{-D^2/\Lambda^2} = K/\pi$, with $K$ the number of Paper~2's
Observation, and the integer-relation search on it, are in Paper~2's
table of examined mechanisms~\cite{paper2}.

% ===================================================================
\section{Two-Loop Identified Constant: $S_{\min}$}
\label{sec:smin}
% ===================================================================

The model's CG-weighted two-loop sunset sum is $2S_{\min}$
(Sec.~\ref{subsec:smin_sunset}), with
\begin{equation}
  S_{\min} = \sum_{k=1}^{\infty} T(k)^2, \qquad
  T(k) = 2\,\zeta(2, k + \tfrac{3}{2})
       - \tfrac{1}{2}\,\zeta(4, k + \tfrac{3}{2}).
  \label{eq:S_min}
\end{equation}
This is a depth-2 multiple Hurwitz zeta value at half-integer shift
$a = 3/2$, a sum over the Dirac spectrum alone; the vertex rule enters
only through the relation $S_{\mathrm{sunset}} = 2S_{\min}$.

\subsection{200-digit computation}

$S_{\min}$ is computed to 200 verified digits using
\texttt{mpmath.nsum} with Levin $u$-transform acceleration.
The first 80 digits agree with three independent methods
(direct summation + explicit tail; Levin $u$-transform; Euler--Maclaurin
26-term asymptotic):
\begin{equation}
  S_{\min} = 2.47993\,69380\,34222\,55441\,35795\,00082\,93821\,44688\ldots
  \label{eq:S_min_value}
\end{equation}

\subsection{Motivic identification: Hoffman double $t$-values}
\label{subsec:smin_identified}

$S_{\min}$ decomposes \emph{exactly} into Hoffman double $t$-values
--- classical level-2 alternating MZVs, periods of $\MT(\Z[1/2])$.
Notation: $\lambda(s) = (1-2^{-s})\zeta(s) = \sum_{o\ge1,\,\mathrm{odd}} o^{-s}$;
$t(b,c) = \sum_{o_2 > o_1 \ge 1,\,\mathrm{odd}} o_2^{-b} o_1^{-c}$
(Hoffman~\cite{hoffman2019odd} descending convention: first index $b$
on the larger summation variable).

\textbf{8-$t$-value decomposition} (verified to residual $1.8\times10^{-129}$ at
130~dps, re-verified at 260~dps):
\begin{align}
  S_{\min} &= 64\bigl[t(2,1) - t(4,1) - 3\,t(2,2) + 3\,t(4,2)
              - t(2,3) + t(4,3) + 3\,t(2,4) - 3\,t(4,4)\bigr]
              \nonumber\\
           &\quad + 32 \sum_{s=3}^{8} d_s\,\bigl(\lambda(s) - 1 - 3^{-s}\bigr),
  \label{eq:Smin_8tvalue}
\end{align}
where the diagonal coefficients are
$d_3=1,\,d_4=-3,\,d_5=-2,\,d_6=6,\,d_7=1,\,d_8=-3$.

\textbf{Stuffle reductions} (exact algebraic identities from the ordering
decomposition $\lambda(b)\lambda(c) = t(b,c)+t(c,b)+\lambda(b+c)$):
\begin{align}
  t(2,2) &= \tfrac{1}{2}\bigl(\lambda(2)^2 - \lambda(4)\bigr), \label{eq:stuffle22}\\
  t(4,4) &= \tfrac{1}{2}\bigl(\lambda(4)^2 - \lambda(8)\bigr), \label{eq:stuffle44}\\
  t(2,4)+t(4,2) &= \lambda(2)\lambda(4) - \lambda(6). \label{eq:stuffle2442}
\end{align}
(The antisymmetric combination $t(2,4)-t(4,2)$ enters $S_{\min}$ with
coefficient zero, so only the reducible symmetric combination appears.)

\textbf{Odd-weight depth-1 reductions} (verified at $\ge 1\times10^{-130}$;
both are verbatim rows of Hoffman~\cite{hoffman2019odd} Appendix~A):
\begin{align}
  t(2,1) &= -\tfrac{7}{16}\,\zeta(3) + \tfrac{1}{8}\,\pi^2\ln 2,
  \label{eq:t21}\\
  t(2,3) &= -\tfrac{31}{64}\,\zeta(5) + \tfrac{1}{16}\,\pi^2\,\zeta(3).
  \label{eq:t23}
\end{align}

\textbf{Remaining odd-weight components --- identified.}  PSLQ at 220~dps
against the full conjectural level-2 bases (dim~8 at weight~5; 19~product
monomials at weight~7) identifies both remaining components (residuals
$2.8\times10^{-221}$ and $3.6\times10^{-222}$):
\begin{align}
  t(4,1) &= \tfrac{1}{96}\,\pi^4\ln 2 - \tfrac{1}{64}\,\pi^2\zeta(3)
            - \tfrac{31}{64}\,\zeta(5),
  \label{eq:t41}\\
  t(4,3) &= \tfrac{1}{128}\,\pi^4\zeta(3) - \tfrac{5}{128}\,\pi^2\zeta(5)
            - \tfrac{127}{256}\,\zeta(7).
  \label{eq:t43}
\end{align}
Both identified forms are likewise verbatim rows of
Hoffman~\cite{hoffman2019odd} Appendix~A ($t(4,1)$ directly; $t(4,3)$
after converting his $t$-basis) --- the PSLQ runs independently
rediscovered the published table (grounding verified against the
source PDF, 2026-06-11).  Assembling the full decomposition yields the
explicit closed form
\begin{equation}
  S_{\min} = 8\pi^2\ln 2 - \tfrac{2}{3}\pi^4\ln 2 - 3\pi^2\zeta(3)
    + \tfrac{1}{2}\pi^4\zeta(3) - \tfrac{5}{2}\pi^2\zeta(5)
    + \tfrac{\pi^6}{4} - \tfrac{3}{2}\pi^4 - \tfrac{\pi^8}{96},
  \label{eq:smin_closed}
\end{equation}
verified against the Hurwitz-tail anchor at 130~dps (residual
$1.7\times10^{-129}$; \texttt{tests/test\_smin\_decomposition.py}).  Two
structural features are exact:\ $\zeta(7)$ cancels against the
$\lambda(7)$ diagonal term, and all rational constants cancel to zero.
Hence $S_{\min}$ lies in $\mathbb{Q}[\pi^2, \ln 2, \zeta(3), \zeta(5)]$:\
the two-loop constant carries \emph{no irreducible depth-2 content at
all} --- every double $t$-value in its decomposition reduces to
depth~$\le 1$.

\subsection{Inner photon sum: closed-form channel count}
\label{subsec:smin_channel}

The inner photon sum $I(n_1, n_2) = \sum_{q\,\mathrm{allowed}}
W(n_1, n_2, q) \cdot d_q^T / \mu_q^p$ (at $p=1$: $d_q^T/\mu_q = 1$)
counts the total SO(4) CG-weighted vertex channels.  Four structural
properties:
\begin{enumerate}
\item[(i)] $I(0, n_2) = 0$ for all $n_2$ (structural zero: vertex parity
  at $n_{\mathrm{ext}}=0$ forces $2n_2$ odd, impossible).
\item[(ii)] $I(n_1, n_2) = I(n_2, n_1)$ (symmetric).
\item[(iii)] $I(n, n) = 2n$.
\item[(iv)] \textbf{Closed form}:
  \begin{equation}
    I(n_1, n_2) = 2\min(n_1, n_2), \qquad n_1, n_2 \ge 0.
    \label{eq:channel_closed_form}
  \end{equation}
  \emph{Proof.}  Take $n_1 \le n_2$ and $d = n_2 - n_1$.  The window
  of Eq.~\eqref{eq:vertex_rule} is $q \in \{d+1, d+3, \ldots,
  n_1 + n_2 - 1\}$, which has $n_1$ elements.  For the component
  $((q+1)/2, (q-1)/2)$ the left triangle
  $((n_1+1)/2, (q+1)/2, (n_2+1)/2)$ holds for $d-1 \le q \le n_1+n_2+1$
  and the right triangle $(n_1/2, (q-1)/2, n_2/2)$ for
  $d+1 \le q \le n_1+n_2+1$; for $((q-1)/2, (q+1)/2)$ the ranges are
  $d+1 \le q \le n_1+n_2+3$ and $d-1 \le q \le n_1+n_2-1$.  All four
  integrality conditions are $n_1 + n_2 + q$ odd.  Each component is
  therefore open on every window $q$, so $W = 2$ there and
  $I = 2n_1$.  \qed\ (Checked against \texttt{geovac/qed\_vertex.py}
  for all $n_1, n_2 \le 8$ in \texttt{tests/test\_smin\_decomposition.py}.)
\end{enumerate}

\subsection{Exact sunset relation}
\label{subsec:smin_sunset}

Substituting $I = 2\min$ into the double sum gives
\begin{equation}
  S_{\mathrm{sunset}} = 2\,S_{\min},
  \label{eq:sunset_relation}
\end{equation}
where $S_{\mathrm{sunset}} = \sum_{n_1,n_2 \ge 1} I(n_1,n_2)\,\phi(n_1)\,\phi(n_2)$
is the CG-exact two-loop sunset.  Its value
$S_{\mathrm{sunset}} = 4.95987387606844510882\ldots$ is verified at
35~dps, and the identity exactly in rational arithmetic at truncation
$n \le 30$, in \texttt{tests/test\_smin\_decomposition.py}.  The two-loop
sunset therefore has the closed form $2\times$Eq.~\eqref{eq:smin_closed}
and carries no irreducible depth-2 content.

% ===================================================================
\section{Three-Loop Results}
\label{sec:three_loop}
% ===================================================================

\subsection{Factorized $O(N^3)$ algorithm}

The three-loop iterated sunset has chain topology: three electron
lines $(n_1, n_2, n_3)$ connected by two photon lines $(q_1, q_2)$,
with vertex selection at each junction.  The chain topology permits
factorization into partial sums:
\begin{equation}
  S^{(3)} = \sum_{n_2=0}^{N}
    \frac{g_{n_2}}{|\lambda_{n_2}|^a}
    \cdot L(n_2) \cdot R(n_2),
  \label{eq:three_loop_factorized}
\end{equation}
where
\begin{equation}
  L(n_2) = \sum_{\substack{n_1, q_1 \\ \text{allowed}}}
    \frac{W(n_1, n_2, q_1) \; g_{n_1} \; d_{q_1}^T}
         {|\lambda_{n_1}|^a \; \mu_{q_1}^p}, \qquad
  R(n_2) = \sum_{\substack{n_3, q_2 \\ \text{allowed}}}
    \frac{W(n_2, n_3, q_2) \; g_{n_3} \; d_{q_2}^T}
         {|\lambda_{n_3}|^a \; \mu_{q_2}^p}.
  \label{eq:LR_partial}
\end{equation}
$W(n_1, n_2, q)$ is symmetric in $(n_1, n_2)$, so $L = R$; the code
computes them independently as a check.  This reduces the five-fold sum
$O(N^5)$ to $O(N^3)$ (each of $L$, $R$ is $O(N^2)$ precomputed,
then the outer sum is $O(N)$).

\noindent\textbf{Verification:} The factorized algorithm matches the
direct $O(N^5)$ sum to $< 10^{-78}$ at $n_{\max} = 10, 15$ across
four parameter combinations (default $a=4,p=1$; vertex-restricted;
higher exponents $a=6,p=2$; with and without CG weights).  The
factorized algorithm at $n_{\max} = 50$ completes in $< 60$ seconds.
Tests: 7 in \texttt{TestFactorizedMatchesDirect} and
\texttt{TestFactorizedConvergence} (all passing; the frozen
factorized-vs-direct gate is $10^{-25}$ --- the $< 10^{-78}$
agreement quoted above is the measured value at the working
precision, re-confirmed 2026-07-04).

\subsection{Convergence analysis}

At default exponents ($a = 4$, $p = 1$) with CG weights, the raw
partial sums converge far too slowly for direct evaluation.  An
empirical power-law fit of the increments over $N \in [20, 120]$ gives
\begin{equation}
  \Delta S^{(3)}(N) \approx 21.6 \cdot N^{-1.23},
  \label{eq:convergence_rate}
\end{equation}
but the fitted exponent is an
artifact: by the channel-count closed
form~\eqref{eq:channel_closed_form} the inner partial sums grow as
$L(n_2) = 4\ln n_2 + O(1)$, so the increments follow
$(\ln N)^2/N^2$ --- a log-modulated law whose local log-slope
rises from $0.9$ to $1.4$ over the fit window and
masquerades as a power law.  The true remaining tail at $N = 120$ is
$\approx 12$ units, i.e.\ \emph{zero} reliable digits from raw partial
sums (cf.\ Table~\ref{tab:three_loop_convergence}: the $n_{\max} = 50$
partial sum is $21.58$ against a limit of $40.66$).

The reliable direct evaluation is a rigorous bracket with no series
acceleration anywhere: the factorized sum has positive terms, so
partial sums with the closed-form full-range inner sum give a strict
lower bound, and the explicit inequalities $L(m) \le 4(\ln m + 1)$,
$\phi(m) \le 2/(m+1)^2$ give an integral tail bound.  At
$N = 4\times10^{6}$,
\begin{equation}
  S^{(3)} \in [40.65339,\ 40.65578]
  \label{eq:s3_bracket}
\end{equation}
(\texttt{tests/test\_s3\_decomposition.py} re-gates it at
$N = 2\times10^{4}$).  The decomposition assembly
(Sec.~\ref{subsec:s3_closure}) gives the canonical value, inside the
bracket:
\begin{equation}
  S^{(3)} = 40.65538\,75635\,85422\,54774\,43399\,69117\ldots
  \label{eq:s3_value}
\end{equation}

The unrestricted three-loop sum factorizes exactly as $D(a)^3$: at
$a = 4$, $D(4)^3 = (\pi^2 - \pi^4/12)^3 \approx 5.379$, purely
$\pi^{\mathrm{even}}$.

\subsection{Decomposition closure: the four trailing-$1$
identifications}
\label{subsec:s3_closure}

Notation: $t_3(b_1,b_2,b_3) = \sum_{o_1>o_2>o_3\ge1\,\mathrm{odd}}
o_1^{-b_1} o_2^{-b_2} o_3^{-b_3}$, the Hoffman triple $t$-value in the
same descending convention as $t(b,c)$.  The channel-count closed
form~\eqref{eq:channel_closed_form} gives
$S^{(3)} = C - 12\,P\,S_{\min} - 9P^3$, with $P = \sum_{n\ge1}\phi(n)
= \pi^2 - \pi^4/12 - 64/81$ and $C$ the three-chain core over odd
$o = 2n+3$; $C$ expands into 12 triple $t$-values, 18 double
$t$-values, and $\lambda$/boundary terms (exact in Fraction arithmetic
at every truncation; \texttt{tests/test\_s3\_decomposition.py}).

The components with trailing argument $b_3 = 1$ carry an odd-harmonic
factor $\widetilde{H}(j) \sim \tfrac12\ln j$, and Levin-type series
acceleration mis-converges on such log-modulated summands with an error
that depends on the working precision.  They are therefore evaluated by
Abel summation over the harmonic weight,
$t_3(b_1,b_2,1) = \sum_{k\ge1} (2k-1)^{-1} A(k)$ with
$A(k) = \sum_{j \ge k} a_j$, where both sums are log-free; with these
values the stuffle witnesses close at $5\times10^{-216}$ and
$8\times10^{-217}$.  The bracket~\eqref{eq:s3_bracket}, which uses no
acceleration, guards against the same failure mode.

PSLQ at 220/340~dps against the complete conjectural level-2
weight-homogeneous bases identifies the four components (residuals $2.6\times10^{-216}$ to $7.3\times10^{-339}$
in the archived dossier run; the frozen regression
\texttt{tests/test\_s3\_decomposition.py} re-gates the identifications
at $10^{-30}$/45~dps):
\begin{align}
  t_3(2,1,1) &= \tfrac{1}{2}\mathrm{Li}_4(\tfrac12)
   + \tfrac{1}{48}\ln^4 2 + \tfrac{1}{24}\pi^2\ln^2 2
   - \tfrac{19}{5760}\pi^4,
  \label{eq:t3_211}\\
  t_3(4,1,1) &= \tfrac{1}{192}\pi^4\ln^2 2
   - \tfrac{1}{64}\pi^2\ln 2\,\zeta(3)
   - \tfrac{11}{322560}\pi^6 - \tfrac{1}{2}\,t(5,1)
   - \tfrac{7}{128}\zeta(3)^2,
  \label{eq:t3_411}\\
  t_3(3,2,1) &= \tfrac{3}{64}\pi^2\ln 2\,\zeta(3)
   - \tfrac{1}{9216}\pi^6
   - \tfrac{1}{2}\,t(5,1) - \tfrac{49}{256}\zeta(3)^2,
  \label{eq:t3_321}\\
  t_3(3,4,1) &= \tfrac{1}{768}\pi^4\ln 2\,\zeta(3)
   + \tfrac{5}{128}\pi^2\ln 2\,\zeta(5)
   + \tfrac{61}{2903040}\pi^8 - \tfrac{1}{256}\pi^2\zeta(3)^2
   \nonumber\\
   &\quad + \tfrac{1}{2}\,t(5,3) - \tfrac{1}{2}\,t(7,1)
   - \tfrac{217}{512}\zeta(3)\zeta(5).
  \label{eq:t3_341}
\end{align}
Every component of the $S^{(3)}$ decomposition except the weight-10
graded piece ($W_{10} = -2048\,t_3(4,3,3) - 1024\,t_3(4,4,2)
- 512\,t_2(8,2) - 1536\,t_2(4,6)$, treated below)
is an identified element of the level-2 ring, with the genuine
depth-2 generators $t(5,1)$ (weight 6) and $t(5,3)$, $t(7,1)$
(weight 8) entering at coefficients $\pm\tfrac12$.
\textbf{Realized motivic depth at three loops: $\le 2$ on every
identified component} --- one rung above $S_{\min}$'s depth $\le 1$,
all catalogued, nothing new minted.  Falsifier:
\texttt{tests/test\_s3\_decomposition.py} (including the
rigorous-bracket regression guard for the Levin/log failure mode).

\paragraph{Full assembly and the weight-10 block.}
Assembling the complete decomposition with all identified components
yields a
stronger statement than the entry-level reductions: \emph{every
depth-2 generator cancels in the total}.  $t(5,1)$, $t(5,3)$,
$t(7,1)$, and $\mathrm{Li}_4(\tfrac12)$ all carry total coefficient
zero, as do $\zeta(7)$, $\zeta(9)$, $\zeta(11)$ and every rational
constant --- the $S^{(3)}$ echo of
$S_{\min}$'s exact $\zeta(7)$ cancellation.  The identified part of
$S^{(3)}$ is a polynomial in $\pi^2$, $\ln 2$, $\zeta(3)$ and
$\zeta(5)$ alone.  The remaining weight-10 block reduces, by exact Fraction
Gaussian elimination on the complete weight-10 quasi-shuffle closure
(rank 157 on 255 words; identity verified numerically at
$7\times10^{-59}$), to
\begin{align}
  W_{10}/256 &= 4\lambda(10) - \lambda(2)\lambda(8)
   - 3\lambda(4)\lambda(6) - 4\lambda(3)\,a(4,3)
   - 2\lambda(4)\,a(4,2) \nonumber\\
  &\quad + a(8,2) + 4\,a(7,3) - 3\,a(6,4)
   - 4\,t_3(2,4,4) - 8\,t_3(3,3,4),
  \label{eq:w10_reduction}
\end{align}
where $a(b,c) := t(b,c) - t(c,b)$ and $\lambda(s) = (1-2^{-s})\zeta(s)$;
the lower-weight $a(4,2)$, $a(4,3)$ are identified objects, leaving
\emph{five} surviving generators: the antisymmetric doubles $a(8,2)$,
$a(7,3)$, $a(6,4)$ and the triples $t_3(2,4,4)$, $t_3(3,3,4)$.  A
structural obstruction protects this block from stuffle methods: the
depth-$\le 3$ relation content of the quasi-shuffle closure saturates
already at the $\lambda\cdot t$ product stage, so no stuffle-product
engineering at any depth can reduce it further.

\paragraph{The weight-10 block closed (2026-06-12).}
The missing input is the Charlton--Hoffman symmetry
theorem~\cite{charltonhoffman2022} --- level-2
regularized-double-shuffle content invisible to the stuffle-only
system.  Its generating-series form (their Thm.~2.21) at depth $m=3$,
$\underline{\phi}=\underline{0}$, extracted at the coefficient of
$y_1^{a-1}y_2^{b-1}y_3^{c-1}$ by exact-Fraction series algebra, gives
the reversal-pair sums with every product term explicit (each instance
used was verified numerically at $10^{-78}$ or better;
\texttt{tests/test\_s3\_w10\_identification.py}):
\begin{align}
  t_3(4,4,2) + t_3(2,4,4) &= \tfrac{5}{6144}\pi^4\zeta(3)^2
   - \tfrac{5}{256}\pi^2\zeta(3)\zeta(5)
   - \tfrac{1}{256}\pi^2\zeta(5,3)
   + \tfrac{137}{99532800}\pi^{10},
  \label{eq:ch_pair1}\\
  t_3(4,3,3) + t_3(3,3,4) &= \tfrac{37}{12288}\pi^4\zeta(3)^2
   - \tfrac{25}{1024}\pi^2\zeta(3)\zeta(5)
   - \tfrac{889}{2048}\zeta(3)\zeta(7)
   - \tfrac{1}{256}\pi^2\zeta(5,3)
   + \tfrac{251}{58060800}\pi^{10},
  \label{eq:ch_pair2}
\end{align}
where $\zeta(5,3) = \sum_{m>n\ge 1} m^{-5}n^{-3}$ is the classical
weight-8 depth-2 MZV generator (conjecturally irreducible over
products of single zetas), entering through the weight-8 layer of the
theorem's $\zeta$-side product terms.  The pair \emph{differences}
follow from four ordinary stuffles $t(x)\,t_2(y,z)$,
$(x,y,z) \in \{(2,4,4),(4,2,4),(3,3,4),(3,4,3)\}$, whose pairwise
combinations eliminate the auxiliary triples $t_3(4,2,4)$,
$t_3(3,4,3)$.  Solving the two $2\times2$ systems identifies all four
weight-10 triples individually:
\begin{align}
  t_3(2,4,4) &= \tfrac{1}{4}a(8,2) - \tfrac{1}{4}a(6,4)
   - \tfrac{1}{6144}\pi^4\zeta(3)^2
   - \tfrac{5}{512}\pi^2\zeta(3)\zeta(5)
   - \tfrac{1}{512}\pi^2\zeta(5,3)
   + \tfrac{373}{174182400}\pi^{10},
  \label{eq:t3_244_closed}\\
  t_3(3,3,4) &= \tfrac{1}{4}a(7,3) - \tfrac{1}{4}a(6,4)
   + \tfrac{1}{12288}\pi^4\zeta(3)^2
   + \tfrac{5}{1024}\pi^2\zeta(3)\zeta(5)
   - \tfrac{889}{4096}\zeta(3)\zeta(7)
   - \tfrac{1}{512}\pi^2\zeta(5,3)
   + \tfrac{251}{116121600}\pi^{10},
  \label{eq:t3_334_closed}
\end{align}
each gated against its independently computed 200-digit value at
residual $\le 6\times10^{-197}$ (likewise $t_3(4,4,2)$,
$t_3(4,3,3)$).

\paragraph{Collapse theorem and the full closed form.}
Substituting Eqs.~\eqref{eq:t3_244_closed}--\eqref{eq:t3_334_closed}
into the full $S^{(3)}$ assembly, the three antisymmetric doubles
$a(8,2)$, $a(7,3)$, $a(6,4)$ \emph{cancel identically}: the explicit
$a$-content of the weight-10 corollary is exactly absorbed by the
$a$-content of the two triples (symbolic verification in the frozen
falsifier).  The weight-10 block therefore contributes
\begin{equation}
  S^{(3)} - \mathrm{Ident} \;=\;
  -7\pi^4\zeta(3)^2 + 70\pi^2\zeta(3)\zeta(5)
  + 6\pi^2\,\zeta(5,3) + \tfrac{179}{170100}\pi^{10},
  \label{eq:s3_w10_contribution}
\end{equation}
where $\mathrm{Ident}$ is the identified part of the symbolic
assembly, and the complete closed form
\begin{equation}
  S^{(3)} \;\in\;
  \mathbb{Q}\!\left[\pi^2,\ \ln 2,\ \zeta(3),\ \zeta(5),\
  \zeta(5,3)\right]
  \label{eq:s3_ring}
\end{equation}
is explicit:
\begin{align}
  S^{(3)} &= 128\pi^2\ln^2 2 - \tfrac{32}{3}\pi^4\ln^2 2
   + 64\pi^2\ln 2 - \tfrac{304}{3}\pi^4\ln 2 + 16\pi^6\ln 2
   - \tfrac{2}{3}\pi^8\ln 2 \nonumber\\
  &\quad - 96\pi^2\ln 2\,\zeta(3) + 16\pi^4\ln 2\,\zeta(3)
   - 80\pi^2\ln 2\,\zeta(5)
   + 20\pi^2\zeta(3)^2 - 7\pi^4\zeta(3)^2 + 70\pi^2\zeta(3)\zeta(5)
   \nonumber\\
  &\quad - 24\pi^2\zeta(3) + 40\pi^4\zeta(3) - 9\pi^6\zeta(3)
   + \tfrac12\pi^8\zeta(3)
   - 20\pi^2\zeta(5) + 30\pi^4\zeta(5) - \tfrac52\pi^6\zeta(5)
   \nonumber\\
  &\quad + 6\pi^2\,\zeta(5,3)
   + \tfrac{8}{3}\pi^4 + \tfrac{374}{45}\pi^6 - \tfrac{8231}{3780}\pi^8
   + \tfrac{124571}{680400}\pi^{10} - \tfrac{1}{192}\pi^{12},
  \label{eq:s3_closed_form}
\end{align}
with $\zeta(5,3)$ the \emph{only} constant beyond the
depth-$\le 1$ ring, at coefficient $6\pi^2$.  The closed form
reproduces the canonical value~\eqref{eq:s3_value}; an independent
route --- the sum built on the earlier channel count, whose closed form
was gated at 200~dps, plus the correction
$2G + 4P S_{\min} - P^3 - 8PQ - R_3$
($G = \sum \min(o_1,o_2)\,\phi(n_1)\,\phi(n_2)^2$ with $o_i = 2n_i + 3$,
$Q = \sum_n \phi(n)^2$, $R_3 = \sum_n \phi(n)^3$, all over $n \ge 1$)
summed directly without the $t$-value table --- agrees to
$3\times10^{-120}$ at 120~dps.  Two structural readings: (i) the
transient-generator cancellation pattern of the symbolic assembly
($t(5,1)$, $t(5,3)$, $t(7,1)$, $\mathrm{Li}_4(\tfrac12)$ all cancel)
extends through weight 10 --- the level-2 antisymmetric doubles
$a(8,2)$, $a(7,3)$, $a(6,4)$ also carry total coefficient zero, and
the \emph{value} of $S^{(3)}$ is independent of these still-unreduced
generators; (ii) the genuinely new depth-2 content at three loops
enters on the \emph{level-1} side ($\zeta(5,3) \in \MT(\Z)$), while
the level-2-specific content (the $\ln 2$ sector) stays at realized
depth $\le 1$ --- new depth arrives from level 1, new level stays at
depth $\le 1$.  Falsifiers:
\texttt{tests/test\_s3\_w10\_identification.py} (symmetry-pair sums,
individual triples, symbolic $a$-collapse) and
\texttt{tests/test\_s3\_decomposition.py} (closed
form~\eqref{eq:s3_closed_form} and the independent route, with
$\zeta(5,3)$ computed live).

\paragraph{Literature grounding (verified against the sources).}
All weight-$\le 7$ closed forms in this section and
Sec.~\ref{subsec:smin_identified} are verbatim rows of Hoffman's
Appendix~A~\cite{hoffman2019odd} --- including the triples
$t_3(2,1,1)$, $t_3(4,1,1)$, $t_3(3,2,1)$, after eliminating
$\zeta(\bar 5,1)$ via his own $t(5,1)$ row.  The PSLQ identifications
above therefore independently rediscovered the published table --- a
two-route verification.  The weight-$\ge 8$ reductions rest on the
theorem-grade chain: Hoffman Cor.~4.1 (exact conversion to alternating
MZVs) $+$ parity/reduction theorems for alternating
sums~\cite{borwein1995,panzer2017} $+$ the data mine's proven
weight-$\le 12$ reductions~\cite{bbv2010}.  That $t_3(3,4,1)$ reduces
to depth $\le 2$ is an instance of the depth-3 parity result for
multiple $t$-values of Wang and Xu~\cite{wangxu2025} (their Cor.~4.3:
at weight $8$ and depth $3$ the parity combination is $2\,t_3$, which
reduces to depth $\le 2$), so only the explicit coefficients of
Eq.~\eqref{eq:t3_341} are ours.  We have not found those coefficients
stated in~\cite{hoffman2019odd}, in the explicit families
of~\cite{charltonhoffman2022} (whose $t$-value convention reverses the
argument order), or in~\cite{wangxu2025}, whose only explicit depth-3
evaluations are at weight $\le 6$.

% ===================================================================
\section{Depth-$k$ Tower}
\label{sec:depth_k}
% ===================================================================

\begin{proposition}[Depth-$k$ tower, corrected]
\label{prop:depth_k}
At $k$ loops on $S^3$ with vertex restrictions and CG weights:
\begin{itemize}
\item $k = 1$: free (unrestricted) integer-$s$ traces of $D^2$
      (equivalently $D(s)$ at even $s$) give $\pi^{\mathrm{even}}$ only
      (Theorem~\ref{thm:T9}, proven);\ vertex-restricted quantities
      need not (see $k = 2$).
\item $k = 2$ (vertex-parity sums):\ the two-loop sunset vertex
      restriction brings $G$ and $\beta(4)$ ($\beta(4)$ first enters here among the identified quantities;\ the $\chi_{-4}$ theorem).
\item $k = 2$ (CG-weighted):\ $S_{\mathrm{sunset}} = 2S_{\min}$ is an identified element of $\MT(\Z[1/2])$
      at level~2.  The depth is \emph{bounded} by 2 (the nesting depth
      of the min-weight double sum) but the content decomposes into
      Hoffman double $t$-values, which are classical period-ring
      elements (Sec.~\ref{subsec:smin_identified}).
\item $k = 3$: $S^{(3)}$ is a fully identified element of
      $\mathbb{Q}[\pi^2, \ln 2, \zeta(3), \zeta(5), \zeta(5,3)]$ with
      \emph{realized} depth $\le 2$ unconditionally
      (Eq.~\eqref{eq:s3_ring}): the sole depth-2 constant is the
      level-1 generator $\zeta(5,3)$, at coefficient $6\pi^2$; every
      level-2 depth-2 generator ($t(5,1)$, $t(5,3)$, $t(7,1)$,
      $a(8,2)$, $a(7,3)$, $a(6,4)$) appears transiently and cancels in
      the total (Sec.~\ref{subsec:s3_closure}).
\item $k = 4$: the \emph{realized} depth of $S^{(4)}$ is $\le 3$.
      Depth-4 content enters $S^{(4)}$ only through the four-chain core
      $C_4$ (the o-space relation
      $S^{(4)} = C_4 - 6PC - 3S_o^2 + 27S_oP^2 - 27P^4$, with
      $S_o = \sum \min(o_1,o_2)\,\phi(n_1)\phi(n_2)$, gives $C_4$
      coefficient $+1$; the other cores are $\le$ depth-3;
      \texttt{tests/test\_s4\_stage1.py}), so the verdict reduces to whether $C_4$'s
      depth-4 part reduces to depth $\le 3$. At weights $9, 11, 13$ the
      quasi-shuffle (stuffle) relation system spans the \emph{entire}
      depth-4 word space (rank $=$ number of depth-4 words: $35, 84, 165$;
      exact rank over two independent primes), collapsing $C_4$'s depth-4
      content to depth $\le 3$ unconditionally. At weights $5, 7$ the
      stuffle system is rank-deficient, but these weights lie below the
      threshold for any depth-$\ge 3$ irreducible: weight-5 values reduce
      to single polylog constants (depth $\le 1$) and weight-7 to depth
      $\le 2$ (the $w \le 7$ GeoVac multiple-$t$ values match Hoffman's
      Appendix-A reductions, verified previously). The explicit closed
      form of $S^{(4)}$ remains open; the depth verdict does not depend on
      it.
\end{itemize}
Loop order bounds the depth of the classical level-2 multiple-$t$-value
content; it does not mint new irreducible transcendentals automatically.
Realized depth so far grows strictly more slowly than loop order
($\le k-1$ at $k = 2, 3, 4$).
\end{proposition}

The mechanism: each vertex coupling creates one level of Hurwitz nesting.
One vertex produces a depth-2 double $t$-value sum ($S_{\min}$).
Two vertices produce a depth-$\le 3$ triple $t$-value sum ($S^{(3)}$).
The prediction extends: $k$ vertices produce a depth-$\le (k+1)$
multiple $t$-value in $\MT(\Z[1/2])$.  Whether the surviving
irreducible depth-$k$ component is new (not reducible by parity/stuffle
within the level-2 algebra) is a question for each $k$ separately --- the
structural nesting is genuine, but the motivic content is classical at
$k = 2$ and, on all identified components, at $k = 3$ (realized depth
$\le 2$).  This is the $S^3$ analog of the increasing motivic depth
observed in flat-space multi-loop Feynman
integrals~\cite{laporta2002,bailey2001}.

% ===================================================================
\section{Fine-Structure Euler Sum Cancellation}
\label{sec:euler_cancel}
% ===================================================================

The Dirac Dirichlet series $D(s)$ studied in the preceding sections
sums over the \emph{free} Camporesi--Higuchi spectrum with uniform
degeneracy weights~$g_n$.  A complementary probe arises from the
\emph{exact} Dirac hydrogen spectrum, where the degeneracy-weighted
energy expansion in powers of $(Z\alpha)^2$ generates Dirichlet sums
with state-dependent coefficients that mix harmonic numbers of varying
depth.

\subsection{Setup}

The exact Dirac energy for quantum numbers $(n,j)$ is
\begin{equation}
  E_{nj} = mc^2 \biggl(
    1 + \frac{(Z\alpha)^2}
    {\bigl(n_r + \sqrt{\kappa^2 - (Z\alpha)^2}\bigr)^2}
  \biggr)^{\!-1/2} - mc^2\,,
\end{equation}
where $\kappa = j + \frac{1}{2}$ and $n_r = n - \kappa$.
Expanding in powers of $x = (Z\alpha)^2$,
\begin{equation}
  \sum_j d(n,j)\, \frac{E_{nj}}{mc^2}
  = \sum_{p=1}^{\infty} c_p(n)\, x^p\,,
\end{equation}
where $d(n,j) = 2(2j+1)$ for $j < n - \tfrac{1}{2}$ and
$d(n,j) = 2j+1$ for $j = n - \tfrac{1}{2}$.
The degeneracy-weighted Dirichlet sums are
\begin{equation}
  D_p \;=\; \sum_{n=1}^{\infty} c_p(n)\,.
\end{equation}

\subsection{Closed forms for $c_p(n)$}

Symbolic expansion of the exact Dirac formula and summation over $j$
gives closed forms involving generalized harmonic numbers
$H_{n-1}^{(r)} = \sum_{k=1}^{n-1} k^{-r}$:
\begin{align}
  c_2(n) &= -\frac{5n - 4}{4n^3}\,, \\[3pt]
  c_3(n) &= \frac{19n-20}{8n^5}
    - \frac{3}{2}\,\frac{H_{n-1}}{n^4}
    - \frac{1}{2}\,\frac{H_{n-1}^{(2)}}{n^3}\,.
\end{align}
At order $p = 2$ and $p = 3$, the sums $D_2$ and $D_3$ collapse to
pure single-zeta values:
\begin{equation}
  D_2 = -\frac{5}{4}\,\zeta(2) + \zeta(3)\,,
  \qquad
  D_3 = \frac{19}{8}\,\zeta(4) - \frac{11}{4}\,\zeta(5)\,.
\end{equation}

\subsection{$\zeta(3)$ complementarity}

The $p = 2$ coefficient defines the \emph{fine-structure shell sum}
\begin{equation}
  \mathcal{S}(n) \;=\; \sum_j d(n,j)\Bigl(\frac{n}{j+\frac{1}{2}}
  - \frac{3}{4}\Bigr) \;=\; \frac{n(5n-4)}{2}\,,
\end{equation}
so that $c_2(n) = -\mathcal{S}(n)/(2n^4)$.
It splits as $-3n^2/2 + 2n(2n-1)$, the second piece twice a hexagonal
number.

In the Fock convention ($n \ge 1$), the single-chirality Dirac mode
degeneracy is $g_n^{\mathrm{Dirac}} = 2n(n+1)$.

\begin{theorem}[$\zeta(3)$ complementarity]
\label{thm:zeta3_complement}
For all $n \ge 1$,
\begin{equation}
\label{eq:shell_complement}
  \mathcal{S}(n) + g_n^{\mathrm{Dirac}} \;=\; \frac{9n^2}{2}
  \;=\; |\lambda_0|^2 \times n^2_{\mathrm{Schr}}\,,
\end{equation}
where $|\lambda_0| = 3/2$ is the lowest Dirac eigenvalue on~$S^3$
and $n^2_{\mathrm{Schr}} = 2n^2$ is the Schr\"odinger degeneracy
(including spin).
At the Dirichlet series level,
\begin{equation}
\label{eq:dirichlet_complement}
  D_{\mathcal{S}}(s) + D_{g}(s) \;=\; \frac{9}{2}\,\zeta(s{-}2)
  \qquad \text{for all } s\,,
\end{equation}
where $D_{\mathcal{S}}(s) = \sum_{n \ge 1} \mathcal{S}(n)\,n^{-s}$
and $D_{g}(s) = \sum_{n \ge 1} g_n^{\mathrm{Dirac}}\,n^{-s}$.
The $\zeta(s{-}1)$ terms cancel identically.
\end{theorem}

\begin{proof}
$\mathcal{S}(n) = 5n^2/2 - 2n$ and $g_n^{\mathrm{Dirac}} = 2n^2 + 2n$;
their sum is $9n^2/2$.  Therefore
$D_{\mathcal{S}}(s) = \tfrac{5}{2}\,\zeta(s{-}2) - 2\,\zeta(s{-}1)$
and $D_{g}(s) = 2\,\zeta(s{-}2) + 2\,\zeta(s{-}1)$; summing gives
$\tfrac{9}{2}\,\zeta(s{-}2)$.
\end{proof}

The period content is in the cancelled term: at odd~$s$,
$\zeta(s{-}1)$ is an even zeta value, but at even~$s$ it is an odd
one ($\zeta(3)$, $\zeta(5)$, \ldots).  At $s = 4$:
\begin{equation}
  D_{\mathcal{S}}(4) = \frac{5}{2}\,\frac{\pi^2}{6} - 2\zeta(3),
\end{equation}
so the $-2n$ piece of $\mathcal{S}$ carries $\zeta(3)$, with sign
opposite to the $+2\zeta(3)$ in $D_{g}(4) = 2\zeta(2) + 2\zeta(3)$.
Their sum $\tfrac{9}{2}\zeta(2) = 3\pi^2/4$ lies in $\mathbb{Q}[\pi^2]$.
This is an analogue, in the fine-structure sector, of T9's
integer-$s$ statement (Theorem~\ref{thm:T9}), not T9 itself, which
concerns the half-integer spectrum $n + 3/2$; the complementarity is a
sum over integer $n$ with $g = 2n(n+1)$.

\subsection{$D_4$ at weight 7: product zeta survives}

At order $p = 4$,
\begin{align}
  c_4(n) &= -\frac{5(41n - 40)}{64n^7}
    + \frac{15}{4}\,\frac{H_{n-1}}{n^6}
    - \frac{1}{4}\,\frac{H_{n-1}^{(2)}}{n^5}
    \notag \\
    &\quad - \frac{3}{4}\,\frac{H_{n-1}^{(3)}}{n^4}
    - \frac{1}{4}\,\frac{H_{n-1}^{(4)}}{n^3}\,.
\end{align}
The harmonic-number sums decompose into weight-7 Euler sums
$S_{r,s} = \sum_{n \ge 1} H_n^{(r)} / n^s$ via the identity
$H_{n-1}^{(r)} = \zeta(r) - \zeta(r,n)$ and Euler's formula
$S_{1,6} = 4\zeta(7) - \zeta(2)\zeta(5) - \zeta(3)\zeta(4)$~\cite{flajolet1998}.

\begin{observation}[Euler sum cancellation at weight~7]
\label{obs:euler_cancel}
At 200~dps, PSLQ identifies
\begin{equation}
\label{eq:D4_decomp}
  D_4 = -\frac{205}{64}\,\zeta(6)
  + \frac{71}{8}\,\zeta(7)
  - \frac{9}{2}\,\zeta(3)\zeta(4)\,,
\end{equation}
with residual $< 10^{-195}$ (archived PSLQ run; the value is pinned
live in \texttt{tests/test\_paper28\_gap\_closure\_pins.py}, and the
$D_4$/$D_5$ decimals were independently cross-verified by direct
partial summation, 2026-07-04).  The product $\zeta(2)\zeta(5)$ cancels
\emph{algebraically}: its coefficient vanishes exactly (verified in
exact rational arithmetic from the Euler sum decomposition), while
$\zeta(3)\zeta(4)$ survives with coefficient $-9/2$.
\end{observation}

This breaks the pure-single-zeta pattern: $D_2$ and $D_3$ involve
only $\zeta(s)$ values, whereas $D_4$ introduces the first zeta
product.  The four weight-7 Euler sums, independently identified by
PSLQ, are
\begin{center}
\begin{tabular}{lccc}
\hline
Sum & $\zeta(7)$ & $\zeta(2)\zeta(5)$ & $\zeta(3)\zeta(4)$ \\
\hline
$S_{1,6}$ & $4$ & $-1$ & $-1$ \\
$S_{2,5}$ & $-10$ & $5$ & $2$ \\
$S_{3,4}$ & $18$ & $-10$ & $0$ \\
$S_{4,3}$ & $-17$ & $10$ & $1$ \\
$S_{5,2}$ & $11$ & $-4$ & $-2$ \\
\hline
\end{tabular}
\end{center}
The stuffle relations $S_{r,s} + S_{s,r} = \zeta(r)\zeta(s) + \zeta(r+s)$
are satisfied to $< 10^{-200}$.

\subsection{$D_5$ at weight 9: product survival rule}

At order $p = 5$,
\begin{align}
  c_5(n) &= \frac{7(71n - 72)}{128\,n^9}
    - \frac{105}{16}\,\frac{H_{n-1}}{n^8}
    + \frac{45}{16}\,\frac{H_{n-1}^{(2)}}{n^7}
    + \frac{5}{4}\,\frac{H_{n-1}^{(3)}}{n^6}
    \notag \\
    &\quad - \frac{3}{8}\,\frac{H_{n-1}^{(4)}}{n^5}
    - \frac{15}{32}\,\frac{H_{n-1}^{(5)}}{n^4}
    - \frac{5}{32}\,\frac{H_{n-1}^{(6)}}{n^3}\,,
\end{align}
verified as exact rational equality against direct sympy expansion for
$n = 1, \ldots, 15$.
The Dirichlet sum decomposes via weight-9 Euler sums
$S_{r,s}$ for $r + s = 9$ (computed via Hurwitz zeta to 250~dps,
stuffle relations verified to $< 10^{-200}$):
\begin{center}
\begin{tabular}{lcccc}
\hline
Sum & $\zeta(9)$ & $\zeta(2)\zeta(7)$ & $\zeta(3)\zeta(6)$ & $\zeta(4)\zeta(5)$ \\
\hline
$S_{2,7}$ & $-35/2$ & $7$ & $2$ & $4$ \\
$S_{3,6}$ & $85/2$ & $-21$ & $0$ & $-6$ \\
$S_{4,5}$ & $-125/2$ & $35$ & $0$ & $5$ \\
$S_{5,4}$ & $127/2$ & $-35$ & $0$ & $-4$ \\
$S_{6,3}$ & $-83/2$ & $21$ & $1$ & $6$ \\
$S_{7,2}$ & $37/2$ & $-6$ & $-2$ & $-4$ \\
\hline
\end{tabular}
\end{center}
PSLQ at 200~dps identifies
\begin{equation}
\label{eq:D5_decomp}
  D_5 = \frac{497}{128}\,\zeta(8)
  - \frac{467}{16}\,\zeta(9)
  + \frac{385}{32}\,\zeta(3)\zeta(6)
  + \frac{75}{8}\,\zeta(4)\zeta(5)\,,
\end{equation}
with residual $< 10^{-190}$.  The product $\zeta(2)\zeta(7)$ cancels
algebraically (verified from the Euler sum collection), while both
$\zeta(3)\zeta(6)$ and $\zeta(4)\zeta(5)$ survive.  This extends the
pattern from $D_4$: product terms grow from 0 (at $D_2$, $D_3$) to 1
(at $D_4$) to 2 (at $D_5$).

\subsection{$D_6$ at weight 11: product survival rule extended}

At order $p = 6$ the weight-11 Euler sums $S_{r,s}$ for $r + s = 11$
are computed via Hurwitz zeta to 400~dps, with stuffle relations
verified to $< 10^{-300}$.  Seven nontrivial sums ($r = 1, \ldots, 8$,
$r \ne 3$) contribute through the $\beta_r$ coefficients of the
$c_6(n)$ expansion.  Individual decompositions are found by PSLQ
against the weight-11 depth-1 MZV basis
$\{\zeta(11), \zeta(2)\zeta(9), \zeta(3)\zeta(8), \zeta(4)\zeta(7),
\zeta(5)\zeta(6)\}$:
\begin{center}
\begin{tabular}{lccccc}
\hline
Sum & $\zeta(11)$ & $\zeta(2)\zeta(9)$ & $\zeta(3)\zeta(8)$ & $\zeta(4)\zeta(7)$ & $\zeta(5)\zeta(6)$ \\
\hline
$S_{1,10}$ & $6$ & $-1$ & $-1$ & $-1$ & $-1$ \\
$S_{2,9}$ & $-27$ & $9$ & $2$ & $6$ & $4$ \\
$S_{4,7}$ & $-329/2$ & $84$ & $0$ & $21$ & $4$ \\
$S_{5,6}$ & $463/2$ & $-126$ & $0$ & $-21$ & $0$ \\
$S_{6,5}$ & $-461/2$ & $126$ & $0$ & $21$ & $1$ \\
$S_{7,4}$ & $331/2$ & $-84$ & $0$ & $-20$ & $-4$ \\
$S_{8,3}$ & $-82$ & $36$ & $1$ & $15$ & $6$ \\
\hline
\end{tabular}
\end{center}
The $\zeta(3)\zeta(8)$ column has a sparse support pattern: nonzero
only for $S_{1,10}$, $S_{2,9}$, and $S_{8,3}$.

Collecting through the $\beta_r$ coefficients and the rational
$\zeta(10)$, $\zeta(11)$ parts:
\begin{equation}
\label{eq:D6_decomp}
  D_6 = -\frac{2289}{512}\,\zeta(10)
  + \frac{1589}{16}\,\zeta(11)
  - \frac{1617}{64}\,\zeta(3)\zeta(8)
  - \frac{1785}{64}\,\zeta(4)\zeta(7)
  - \frac{2065}{64}\,\zeta(5)\zeta(6)\,,
\end{equation}
verified to 62~digits against direct Euler-sum evaluation at
400~dps.  The product $\zeta(2)\zeta(9)$ cancels algebraically,
and all three surviving products $\zeta(3)\zeta(8)$,
$\zeta(4)\zeta(7)$, $\zeta(5)\zeta(6)$ are present, matching the
prediction $\max(0, \lfloor(2 \cdot 6 - 5)/2\rfloor) = 3$.

\begin{observation}[Product survival rule]
\label{obs:product_rule}
At each Sommerfeld order $p \ge 2$, the Dirichlet sum $D_p$ decomposes
into weight-$(2p{-}1)$ zeta values.  The product $\zeta(2)\zeta(2p{-}3)$
cancels algebraically at every order.  Exactly
$\max(0, \lfloor(2p{-}5)/2\rfloor)$ product terms survive, all of the
form $\zeta(a)\zeta(b)$ with $a + b = 2p{-}1$, $a$ odd $\ge 3$,
$b$ even $\ge 4$.
\end{observation}

Verified through $D_6$ ($p = 6$, weight~11); no derivation of the rule
is known.

\subsection{Transcendental structure}

Substituting $\zeta(6) = \pi^6/945$ and $\zeta(4) = \pi^4/90$,
equation~\eqref{eq:D4_decomp} becomes
\begin{equation}
\label{eq:D4_three_tier}
  D_4 = -\frac{41\pi^6}{12096}
  + \frac{71}{8}\,\zeta(7)
  - \frac{\pi^4}{20}\,\zeta(3)\,.
\end{equation}
This is a three-tier decomposition: $\pi^6$ (weight-6 even-zeta),
$\zeta(7)$ (weight-7 odd-zeta), and the mixed product
$\pi^4 \cdot \zeta(3)$ (weight-4 even $\times$ weight-3 odd).
In Paper~18's operator-order taxonomy, $\pi^{\mathrm{even}}$ comes
from $\Delta_{\mathrm{LB}}$ and $D^2$ (not from every second-order
operator; Theorem~\ref{thm:T9}'s scope) and $\zeta(\mathrm{odd})$ from
first-order operators; the mixed product $\pi^4 \cdot \zeta(3)$ has
both tiers multiplicatively in a single observable.

The $\zeta(2)\zeta(5)$ cancellation is non-trivial: four of the five
contributing Euler sums carry nonzero $\zeta(2)\zeta(5)$ coefficients
($S_{1,6}$, $S_{2,5}$, $S_{4,3}$, $S_{5,2}$), and the cancellation
emerges only in the specific linear combination that the Dirac
degeneracy weights and harmonic-number structure select.  In contrast,
$S_{3,4}$ alone is free of $\zeta(3)\zeta(4)$, but the remaining four
sums contribute $\zeta(3)\zeta(4)$ terms that do not cancel.

% ===================================================================
\section{Comprehensive Transcendental Tables}
\label{sec:tables}
% ===================================================================

\begin{table}[h]
\centering
\caption{Even/odd decomposition of $D(s)$: PSLQ at 80+ dps.}
\label{tab:D_even_odd}
\begin{tabular}{@{}clll@{}}
\toprule
$s$ & $D_{\mathrm{even}}(s)$ & $D_{\mathrm{odd}}(s)$ & $D_{\mathrm{even}} - D_{\mathrm{odd}}$ \\
\midrule
4 & $\frac{\pi^2}{2} - \frac{\pi^4}{24} - 4G + 4\beta(4)$
  & $\frac{\pi^2}{2} - \frac{\pi^4}{24} + 4G - 4\beta(4)$
  & $8(\beta(4) - G)$ \\[4pt]
5 & odd-zeta + Dirichlet
  & odd-zeta $-$ Dirichlet
  & $16(\beta(5) - \beta(3))$ \\[4pt]
6 & $\pi^{\mathrm{even}}$ + Dirichlet
  & $\pi^{\mathrm{even}}$ $-$ Dirichlet
  & $32(\beta(6) - \beta(4))$ \\
\bottomrule
\end{tabular}
\end{table}

\begin{table}[h]
\centering
\caption{Verification of Theorem~\ref{thm:chi4}: $D_{\mathrm{even}}(s) -
D_{\mathrm{odd}}(s) = 2^{s-1}(\beta(s) - \beta(s{-}2))$.}
\label{tab:chi4_verification}
\begin{tabular}{@{}ccc@{}}
\toprule
$s$ & $2^{s-1}$ & Residual $|$LHS $-$ RHS$|$ \\
\midrule
2 & 2 & 0 (sympy) \\
3 & 4 & 0 (sympy) \\
4 & 8 & $< 10^{-100}$ \\
5 & 16 & $< 10^{-100}$ \\
6 & 32 & $< 10^{-100}$ \\
7 & 64 & $< 10^{-100}$ \\
8 & 128 & $< 10^{-100}$ \\
9 & 256 & $< 10^{-100}$ \\
10 & 512 & $< 10^{-100}$ \\
15 & 16384 & $< 10^{-100}$ \\
\bottomrule
\end{tabular}
\end{table}

\begin{table}[h]
\centering
\caption{Three-loop factorized convergence at $a = 4$, $p = 1$,
CG-weighted.  The limit is $S^{(3)} = 40.6554\ldots$
(Eq.~\eqref{eq:s3_value}): raw partial sums carry zero reliable digits
at these truncations --- the $(\ln N)^2/N^2$ increment law, not the
naively fitted power law, governs the tail
(Sec.~\ref{sec:three_loop}).  Rows computed from
\texttt{geovac/qed\_three\_loop.py}, checked against an independent
exact-rational route, and pinned in
\texttt{tests/test\_paper28\_gap\_closure\_pins.py}.}
\label{tab:three_loop_convergence}
\begin{tabular}{@{}cccc@{}}
\toprule
$n_{\max}$ & $S^{(3)}(n_{\max})$ & $\Delta$ from previous & Time (s) \\
\midrule
5  & 2.753 & --- & $< 0.1$ \\
10 & 6.836 & 4.083 & $< 0.1$ \\
15 & 10.181 & 3.345 & 0.2 \\
20 & 12.837 & 2.655 & 0.5 \\
50 & 21.581 & --- & $< 60$ \\
$\infty$ & $40.6554\ldots$ & --- & (decomposition route) \\
\bottomrule
\end{tabular}
\end{table}

% ===================================================================
\section{Discussion}
\label{sec:discussion}
% ===================================================================

The periods met in this paper sort into four layers.  Single spectral
sums of $D^2$ at integer $s$ lie in $\mathbb{Q}[\pi^2]$
(Theorem~\ref{thm:T9}); single sums of $|D|$ at odd $s \ge 5$ are
rational combinations of two odd zeta values (Theorem~\ref{thm:parity});
restricting $n$ to one parity class adds the level-4 values $G$ and
$\beta(4)$, which cancel in the full sum (Theorem~\ref{thm:chi4}).
Nested chain sums lie in the level-2 (alternating) multiple zeta
values with realized depth $\le k-1$ at loop order $k = 2, 3, 4$
(Proposition~\ref{prop:depth_k}).  At two loops every double $t$-value
reduces to depth $\le 1$.  At three loops the closed form
(Eq.~\eqref{eq:s3_closed_form}) has one depth-2 constant, $\zeta(5,3)$,
which is a level-1 period ($\zeta(5,3) \in \MT(\Z)$), while every
level-2 depth-2 generator met along the way cancels: new depth arrives
from level~1, and the $\ln 2$ sector stays at depth $\le 1$.  Sums over
the exact Dirac hydrogen spectrum give weight-$(2p{-}1)$ Euler sums with
a product-cancellation pattern observed through $p = 6$
(Observation~\ref{obs:product_rule}).

The vertex model enters the nested sums only through its window.
Because $W = 2$ on every window triple and $0$ elsewhere
(Sec.~\ref{subsec:smin_channel}), the channel-count weight multiplies
each vertex-restricted sum by a power of $2$, so the period content
above is a property of the window
$|n_1 - n_2| \le q \le n_1 + n_2$, $n_1 + n_2 + q$ odd.  The physical
vector-photon vertex has a different window and nonconstant reduced
amplitudes (Paper~33); the periods of the corresponding sums are not
determined here.

The depth-$k$ tower fits the framework of Zagier's weight
conjecture~\cite{zagier1994}: loop order raises the weight, and the
algebra of multiple zeta values receives the result.  The half-integer
shift $a = 3/2$ of the Dirac spectrum is what puts the sums in the
odd-denominator ($t$-value) form of the level-2 ring.

\subsection{Open questions}

\begin{enumerate}
\item Does the depth-$k$ tower stabilize?  Realized depth is $\le k-1$
      at $k = 2, 3, 4$; the explicit closed form of $S^{(4)}$ and the
      continuation to $k \ge 5$ are open.
\item Is there a Picard--Fuchs interpretation of $S_{\min}$, analogous
      to Ap\'ery's recurrence for $\zeta(3)$~\cite{apery1979}?
\item What are the periods of the chain sums built on the physical
      vector-photon vertex, with its window and its reduced amplitudes
      (Paper~33)?  This needs the normalization of that vertex.
\item Can the product survival rule (Observation~\ref{obs:product_rule})
      be derived for all $p$ from Euler-sum reflection identities?
\item The model's magnetic-moment ratio: its normalization is owed
      (Sec.~\ref{sec:anomalous_moment}).  Within the model, is the
      asymptotic exponent of Eq.~\eqref{eq:mode_power_law} exactly
      $-2$; can the prefactor in $V_{\mathrm{mag}} \to (4/3)/|\lambda|$
      be derived in closed form; and can the telescoping recurrence of
      the exact square-root forms of $V_{\mathrm{mag}}$ and
      $B(n_{\mathrm{int}} = 1)$ (the negative square-root argument at
      level $n$ equals the positive one at level $n - 1$) be proven
      from the coupling rules?
\end{enumerate}

% ===================================================================
% References
% ===================================================================

\begin{thebibliography}{99}

\bibitem{paper18}
GeoVac Project, ``Spectral-Geometric Exchange Constants: Projecting
Discrete Spectra onto Continuous Manifolds,'' Paper~18 (2026).

\bibitem{paper32}
GeoVac Project, ``The GeoVac Spectral Triple: A Synthesis Paper Locking
the Framework into the Marcolli--van~Suijlekom Gauge-Network Lineage,''
Paper~32 (2026).

\bibitem{paper51}
GeoVac Project, ``A Convergent Discrete Replica Derivative on the
Spinor Cone, with the Two-Term Camporesi--Higuchi Spectral Action on
$S^3\times S^1$,'' Paper~51 (2026).

\bibitem{paper33}
GeoVac Project, ``Selection Rules of the Vector-Photon Vertex on $S^3$:
the Physical Vertex, the Vertex Model and the Graph,'' Paper~33 (2026).

\bibitem{paper2}
GeoVac Project, ``The Fine Structure Constant as a Spectral Coincidence
on the Three-Sphere,'' Paper~2 (2025).

\bibitem{parker1980}
L.~Parker and D.~J.~Toms,
\textit{Quantum Field Theory in Curved Spacetime:\ Quantized Fields
and Gravity} (Cambridge University Press, 2009), Chs.~6--7 --- the
curved-spacetime effective-action formalism from which the leading
curvature correction quoted here follows; see also L.~Parker and
D.~J.~Toms, ``Renormalization-group analysis of grand unified
theories in curved spacetime,''
\textit{Phys.~Rev.~D}\ \textbf{29}, 1584 (1984) (background:
curvature-dependent coupling running, not the $g{-}2$ formula).

\bibitem{camporesi1996}
R.~Camporesi and A.~Higuchi, ``On the eigenfunctions of the Dirac
operator on spheres and real hyperbolic spaces,''
\textit{J.~Geom.~Phys.}\ \textbf{20}, 1--18 (1996).

\bibitem{laporta2002}
S.~Laporta, ``High precision $\epsilon$-expansions of massive four-loop
vacuum bubbles,''
\textit{Phys.~Lett.~B}\ \textbf{549}, 115--122 (2002).

\bibitem{zagier1994}
D.~Zagier, ``Values of zeta functions and their applications,''
in \textit{First European Congress of Mathematics}, Vol.~II,
Birkh\"auser, 497--512 (1994).

\bibitem{bailey2001}
D.~H.~Bailey and D.~J.~Broadhurst, ``Parallel integer relation
detection: techniques and applications,''
\textit{Math.~Comp.}\ \textbf{70}, 1719--1736 (2001).

\bibitem{hoffman2019odd}
M.~E.~Hoffman, ``An odd variant of multiple zeta values,''
\textit{Commun.~Number Theory Phys.}\ \textbf{13}, 529--567 (2019);
arXiv:1612.05232.

\bibitem{borwein1995}
D.~Borwein, J.~M.~Borwein, and R.~Girgensohn, ``Explicit evaluation
of Euler sums,''
\textit{Proc.~Edinburgh Math.~Soc.}\ \textbf{38}, 277--294 (1995).

\bibitem{panzer2017}
E.~Panzer, ``The parity theorem for multiple polylogarithms,''
\textit{J.~Number Theory}\ \textbf{172}, 93--113 (2017);
arXiv:1512.04482.

\bibitem{bbv2010}
J.~Bl\"umlein, D.~J.~Broadhurst, and J.~A.~M.~Vermaseren,
``The multiple zeta value data mine,''
\textit{Comput.~Phys.~Commun.}\ \textbf{181}, 582--625 (2010);
arXiv:0907.2557.

\bibitem{charltonhoffman2022}
S.~Charlton and M.~E.~Hoffman, ``Symmetry results for multiple
$t$-values,'' arXiv:2204.14183 (2022).

\bibitem{wangxu2025}
Z.~Wang and C.~Xu, ``Contour integrations and parity results of
cyclotomic Euler $T$-sums and multiple $t$-values,''
arXiv:2509.06706 (2025).

\bibitem{flajolet1998}
P.~Flajolet and B.~Salvy, ``Euler sums and contour integral
representations,''
\textit{J.~Exp.~Math.}\ \textbf{7}, 15--35 (1998).

\bibitem{apery1979}
R.~Ap\'ery, ``Irrationalit\'e de $\zeta(2)$ et $\zeta(3)$,''
\textit{Ast\'erisque}\ \textbf{61}, 11--13 (1979).

\bibitem{rivoal2000}
T.~Rivoal, ``La fonction z\^eta de Riemann prend une infinit\'e de
valeurs irrationnelles aux entiers impairs,''
\textit{C.~R.~Acad.~Sci.~Paris, S\'er.~I}\ \textbf{331}, 267--270
(2000).

\bibitem{chamseddine_connes_uncanny}
A.~H.~Chamseddine and A.~Connes, ``The uncanny precision of the spectral action,''
\textit{Commun.\ Math.\ Phys.}\ \textbf{293}, 867--897 (2010);\ arXiv:0812.0165.

\end{thebibliography}

\end{document}
